%
% notification.tex
%
% Z Specification for Biff Notification System
% Formal model of NATS connection, MCP server, filesystem, and
% status bar states and transitions in the send and receive flow.
%
% Corrected model: eager session capture during server lifespan
% ensures the suspenders path always has a valid session reference.
%

\documentclass[a4paper,10pt,fleqn]{article}
\usepackage[margin=1in]{geometry}
\usepackage{fuzz}
\usepackage[colorlinks=true,linkcolor=blue,citecolor=blue,urlcolor=blue]{hyperref}

\begin{document}

\title{Biff Notification System: A Z Specification}
\author{Formal Model of Push Notification State Transitions}
\date{March 2026}
\hypersetup{
  pdftitle={Biff Notification System: A Z Specification},
  pdfauthor={punt-labs},
  pdfsubject={Z Specification},
  pdfcreator={fuzz/probcli}
}
\maketitle

\tableofcontents
\newpage

%%----------------------------------------------------------------------
\section{Introduction}
%%----------------------------------------------------------------------

Biff's push notification system delivers real-time updates (wall
broadcasts and talk messages) to Claude~Code's status bar through two
independent channels: the MCP protocol (fast, tool-list-changed
notification) and a filesystem status file (slow, polled by the status
bar command).

The notification pipeline has a ``belt and suspenders'' design: tool
handlers fire notifications via the MCP request context (belt path),
while a background poller fires them via a stored
\texttt{ServerSession} reference (suspenders path).  NATS subscription
callbacks and KV watchers only wake the poller from napping; the
poller detects the changes on its next active tick and fires the
notification from the MCP server's own coroutine context.

The session reference is captured eagerly during the server's lifespan
startup, before the NATS connection is established and before the
background poller begins---this is what closes the biff-8g0 defect
class, in which the poller could wake before any session had ever been
captured.  NATS callbacks and KV watchers do not fire notifications
directly---they are unreliable from background coroutine
contexts---but they wake the poller from napping, so the next active
tick (2~seconds) detects the changes via its own reads and attempts
the notification.

\emph{Whether that attempt succeeds is a separate question from
whether a session was ever captured, and Sections~\ref{sec:state}
and~\ref{sec:sessionloss} correct an over-claim in earlier drafts on
exactly this point.}  A send on the stored session reference can fail
independently of NATS connectivity---the client's own MCP transport
can die while the NATS connection underneath stays healthy---and
root cause~\#1 (``dropped \texttt{tools/list\_changed} mid-session, no
recovery,'' Section~\ref{sec:sessionloss}) is what happens when that
failure is not recorded: the session reference is cleared, the
description-changed flag is left set, and nothing tells the next
client reconnect that a notification is owed.  The governing property
is not ``the poller always has a valid session''---that is false, once
a send can fail---but the accounting safety property $notifyLost =
zfalse$: every suspenders failure is recorded, so recovery is
guaranteed rather than accidental.

This specification models the four state spaces---NATS connection,
MCP server session, filesystem, and status bar---and the operations
that transition them during the wall and talk notification flows.

A second delivery mechanism sits alongside the status bar: the server
surfaces incoming activity by mutating a \emph{tool description} and
firing \texttt{tools/list\_changed}.  The \texttt{read\_messages}
description gains \texttt{(N unread)}; the \texttt{talk} description gains
a \texttt{[TALK] \ldots} banner with an accept hint.
Sections~\ref{sec:marker} to~\ref{sec:markerfix2} model these two
channels and establish the governing invariant: for each channel the
marker is present exactly when its activity exists, and a talk marker
always carries a session-naming accept hint.  The read channel is
self-clearing and serves as the reference; the talk channel, as
originally built (biff-9la), was not.  The stale-marker defect is
reproduced here as a reachable ProB state and then proved unreachable
once the reconciling transitions of Section~\ref{sec:markerfix} are in
place.

A third defect sits underneath both marker channels rather than
alongside them: the marker can be server-side correct---mutated,
consistent with $MarkerConsistent$---while the notification that would
tell Claude~Code to re-read it is silently lost.
Section~\ref{sec:sessionloss} models this as root cause~\#1, reproduces
it as a reachable ProB state, and proves the corresponding safety
property once the suspenders failure is properly recorded.  A second
round of review found the same failure mode at the other two sites
that write $pendingNotify$---the capture/reconnect flush
(Section~\ref{sec:capflush}) and the belt path itself
(Section~\ref{sec:beltpath})---and, sharper still, found that treating
each site's read-then-await-then-write as atomic was itself the
coarseness: Section~\ref{sec:race} splits all three into a read/decide
step and a write step, reproduces the resulting lost-update race as a
reachable state, and proves it closed by a single mutex serialising the
three critical sections.

%%----------------------------------------------------------------------
\section{Basic Types}
%%----------------------------------------------------------------------

\begin{zed}
  [USER, MSG, SOURCEKEY, SESSIONKEY]
\end{zed}

$USER$ identifies a biff user.  $MSG$ is an opaque message body.
$SOURCEKEY$ uniquely identifies a display queue item (e.g.\
\texttt{wall:2026-03-01T12:00:00} or \texttt{talk:eric}).
$SESSIONKEY$ is a composite \texttt{user:tty} key naming one specific
session; it is the same notion as the $SESSIONKEY$ of $talk.tex$, and it
is what an accept must target (\texttt{talk @user:tty}).  A pending talk
invite that retains its inviter's $SESSIONKEY$ can render an accept hint
that names a session; one that has dropped it can only render a bare
\texttt{@user}, which is not a runnable accept command.

%%----------------------------------------------------------------------
\section{Free Types}
%%----------------------------------------------------------------------

\begin{zed}
  ZBOOL ::= ztrue | zfalse
\end{zed}

\begin{zed}
  NatsState ::= natsDisconnected | natsConnected | natsClosed
\end{zed}

The NATS client connection lifecycle: initially disconnected, lazily
connected on first use, and permanently closed on shutdown.

\begin{zed}
  SessionState ::= sessNone | sessCaptured
\end{zed}

The MCP \texttt{ServerSession} reference: either not yet captured
(\texttt{sessNone}) or captured during server lifespan startup
(\texttt{sessCaptured}).  Once captured, the session reference
is never invalidated.

\begin{zed}
  ItemKind ::= itemWall | itemTalk
\end{zed}

Display queue items are either persistent wall announcements or
ephemeral talk messages.

\begin{zed}
  NotifyResult ::= notifySent | notifyDropped
\end{zed}

The outcome of a \texttt{notify\_tool\_list\_changed} attempt.
\emph{Erratum}: earlier drafts claimed $notifyDropped$ was only the
initial default value and that every notification-triggering operation
produces $notifySent$.  That is false once a send can genuinely fail
(Section~\ref{sec:sessionloss}): $notifyDropped$ is a real, reachable
outcome at all three sites---suspenders, belt, and the capture/reconnect
flush.  What holds instead is $notifyLost = zfalse$: every site that
produces $notifyDropped$ also re-arms $pendingNotify$, so a dropped send
is recorded rather than silently absorbed.

\begin{zed}
  CriticalSite ::= siteSuspenders | siteBelt | siteCapture
\end{zed}

The three call sites that write $pendingNotify$ as the consequence of an
awaited send: the background poller (Section~\ref{sec:sessionloss}), the
belt path (Section~\ref{sec:notifycore}), and the capture/reconnect flush
(Sections~\ref{sec:sessioncapture} and~\ref{sec:sessionloss}).
Section~\ref{sec:race} splits each into a read/decide step and a
write step and uses $CriticalSite$ to track which sites currently have
a write outstanding between the two.

%%----------------------------------------------------------------------
\section{Global Constants}
%%----------------------------------------------------------------------

The display queue and the unread counter are bounded.  In the running
system the display rotation queue holds at most twenty items and the
unread count is capped at one hundred.  Pinning those concrete numbers
into the model, however, makes exhaustive checking hopeless: a sequence
of up to twenty source keys over even two carriers, crossed with an
unread count ranging over $0 \upto 100$, yields hundreds of thousands
of reachable states before any interesting property is reached.

Following the convention of the sister specification $talk.tex$, the
bounds are therefore not hard-wired.  Each is the cardinality of an
abstract carrier---$QSLOT$ for queue slots, $USLOT$ for unread slots---so
the bound is a \emph{scoping parameter} ProB shrinks for exhaustive
checking rather than a fixed value.  A third carrier, $ASLOT$, plays the
same role for the pending-invite time-to-live: $maxInviteAge$ is its
cardinality, and a pending invite whose age reaches that bound may be
aged out.  A fourth carrier, $PSLOT$, bounds the pending-invite set
itself: $maxPending$ is its cardinality, the number of pending invites
the recipient will retain at once.  This bound is not decoration.  A
pending invite is keyed by the wire-controlled inviter identity; the
time-to-live is eventual, not a cap, so without a size bound each
distinct forged inviter would add a permanent entry, an
unbounded-growth denial of service (the biff-vr4 analogue).  The bound
caps the set exactly as $maxQueueSize$ caps the message queue, and
$TalkInviteArriveOverflow$ evicts to hold it, mirroring $ReceiveOverflow$
of $talk.tex$.  All four carriers are required non-empty, so each bound
is at least one; at $\# ASLOT = 1$ the clock has a single tick, which is
the smallest model in which a fresh invite (age $0$) and an expired one
(age $1$) are distinct.

\begin{zed}
  [QSLOT, USLOT, ASLOT, PSLOT]
\end{zed}

\begin{axdef}
  maxQueueSize : \nat \\
  maxUnreadCount : \nat \\
  maxInviteAge : \nat \\
  maxPending : \nat
  \where
  QSLOT \neq \emptyset \\
  USLOT \neq \emptyset \\
  ASLOT \neq \emptyset \\
  PSLOT \neq \emptyset \\
  maxQueueSize = \# QSLOT \\
  maxUnreadCount = \# USLOT \\
  maxInviteAge = \# ASLOT \\
  maxPending = \# PSLOT
\end{axdef}

Default descriptions and the idle marker are opaque constants:

\begin{axdef}
  wallDescDefault : MSG \\
  talkDescDefault : MSG \\
  readDescDefault : MSG \\
  idleMarker : MSG
\end{axdef}

$readDescDefault$ is the base \texttt{read\_messages} tool description,
carrying no unread marker.  Its role in the read channel mirrors that of
$talkDescDefault$ in the talk channel and $wallDescDefault$ in the wall
channel.

The intended verification is exhaustive model checking at set size~2,
which supplies two distinct carrier elements throughout.  Modelling the
tool-context brackets (Section~\ref{sec:toolcontext}) roughly doubles the
reachable state count and, by making the belt path live, opens the
receive-flow subspace behind $TalkStart$; at the full queue and unread
bounds the check no longer terminates in a practical time.  The bounded
carriers exist precisely for this: shrinking $QSLOT$, $USLOT$ and $ASLOT$
to one keeps the exploration complete while retaining every distinction
the properties depend on.

\begin{verbatim}
probcli docs/notification.tex -model_check -noltl -coverage \
  -p DEFAULT_SETSIZE 2 \
  -card SOURCEKEY 1 -card SESSIONKEY 1 \
  -card QSLOT 1 -card USLOT 1 -card ASLOT 1 -card PSLOT 1 \
  -p MAX_INITIALISATIONS 16 -p MAX_OPERATIONS 8000
\end{verbatim}

Each carrier may be shrunk independently with \texttt{-card}~$X$~$n$
without disturbing the message and key carriers.  $MAX\_INITIALISATIONS$
must exceed the number of identities $Init$ leaves unconstrained, or ProB
silently drops initial states and reports an incomplete exploration; at
these bounds $Init$ is fully determined, so exploration reports
$INITIALISATION$:$8$ against $SETUP\_CONSTANTS$:$8$ with the bound never
binding, and terminates with ``all open states visited''---a complete
check, not a truncated one.

%%----------------------------------------------------------------------
\section{State}
\label{sec:state}
%%----------------------------------------------------------------------

The system state is flattened into a single schema for ProB
compatibility.  It covers four subsystems: NATS connection state,
MCP session capture, the display queue, the filesystem status file,
the status bar, and the last notification outcome.

\begin{schema}{NotificationState}
  natsConn : NatsState \\
  kvProvisioned : ZBOOL \\
  talkSubscribed : ZBOOL \\
  session : SessionState \\
  toolContext : ZBOOL \\
  wallDesc : MSG \\
  talkDesc : MSG \\
  readDesc : MSG \\
  wallDescChanged : ZBOOL \\
  talkDescChanged : ZBOOL \\
  unreadCount : \nat \\
  talkPending : USER \pfun SESSIONKEY \\
  talkBare : \power USER \\
  talkQueued : \nat \\
  talkConnected : ZBOOL \\
  pendingAge : \nat \\
  queueItems : \seq SOURCEKEY \\
  queueKinds : SOURCEKEY \pfun ItemKind \\
  queueTexts : SOURCEKEY \pfun MSG \\
  currentIndex : \nat \\
  fileExists : ZBOOL \\
  fileUnreadCount : \nat \\
  fileItems : \seq SOURCEKEY \\
  barLine2 : MSG \\
  barStale : ZBOOL \\
  lastNotifyResult : NotifyResult \\
  staleWithdrawHit : ZBOOL \\
  pendingNotify : ZBOOL \\
  notifyLost : ZBOOL \\
  inFlight : \power CriticalSite \\
  beltPnSnapshot : ZBOOL \\
  capPnSnapshot : ZBOOL \\
  raceLost : ZBOOL
  \where
  \dom queueKinds = \ran queueItems \\
  \dom queueTexts = \ran queueItems \\
  \# queueItems \leq maxQueueSize \\
  (queueItems \neq \langle \rangle \implies
  currentIndex < \# queueItems) \\
  (queueItems = \langle \rangle \implies
  currentIndex = 0) \\
  fileUnreadCount \leq maxUnreadCount \\
  unreadCount \leq maxUnreadCount \\
  talkQueued \leq maxQueueSize \\
  \# talkPending + \# talkBare \leq maxPending \\
  pendingAge \leq maxInviteAge \\
  ((\dom talkPending = \emptyset \land talkBare = \emptyset) \implies
  pendingAge = 0) \\
  \dom talkPending \cap talkBare = \emptyset \\
  ((talkDesc \neq talkDescDefault) \iff \\
  \quad (\dom talkPending \neq \emptyset \lor talkBare \neq \emptyset \lor \\
  \quad~ talkQueued > 0 \lor talkConnected = ztrue)) \\
  ((readDesc \neq readDescDefault) \iff unreadCount > 0) \\
  (talkDesc \neq talkDescDefault \implies talkBare = \emptyset) \\
  (talkSubscribed = ztrue \implies
  natsConn = natsConnected) \\
  (kvProvisioned = ztrue \implies
  natsConn = natsConnected) \\
  staleWithdrawHit = zfalse \\
  notifyLost = zfalse \\
  \# inFlight \leq 1 \\
  raceLost = zfalse
\end{schema}

Key invariants: queue metadata maps have the same domain as the
queue sequence range; subscription and KV require a live connection;
the current index is always within bounds; the pending-invite set is
bounded, $\# talkPending + \# talkBare \leq maxPending$, so no stream
of forged inviters can grow it without limit; and $notifyLost = zfalse$,
established in Section~\ref{sec:sessionloss}, which says that a send
failure at any of the three $pendingNotify$-writing sites is never
silently absorbed.  The last two, $\# inFlight \leq 1$ and $raceLost =
zfalse$, are Section~\ref{sec:race}'s: the first is the Z-level model of
a mutex serialising the three sites' critical sections, and the second
says that under it no site's write can ever land on a stale read of
another site's concurrent write.

\paragraph{Erratum (root cause \#1, dropped \texttt{tools/list\_changed}
mid-session).}  Earlier drafts of this section carried a further
invariant, $natsConn = natsConnected \implies session = sessCaptured$,
glossed as ``a connected NATS client implies a captured session,'' with
the accompanying claim in Section~\ref{sec:sessioncapture} that ``once
captured, the session reference is never invalidated.''  That claim is
false of the running system.  $notify\_tool\_list\_changed()$'s
suspenders branch clears the stored session reference to \texttt{None}
when a send raises---a live NATS connection survives the send failure
untouched, but the cached session reference does not---so
$natsConn = natsConnected$ and $session = sessNone$ can co-occur in the
real system.  The invariant as stated is refuted by exactly the
transition that models this failure ($PollTickNotifyFail$,
Section~\ref{sec:sessionloss}); keeping it would make that transition
infeasible and hide the defect rather than expose it.  It is removed
here.  The property it was protecting---that $EnsureConnected$ can only
run after a session has been captured at least once, the biff-8g0
guard---did not actually depend on session capture being permanent; it
depended only on capture happening \emph{before the first connect}.
That is now stated as an explicit precondition on $EnsureConnected$
itself (Section~\ref{sec:ensureconnected}), where it belongs, rather
than as a state invariant that also asserts something stronger and
false.

\paragraph{The tool-description channels.}  Three of the components above
model tool descriptions that Claude~Code shows to the agent, and two more
model the underlying activity that a description marker is supposed to
track.  These are the subject of the correctness argument in
Section~\ref{sec:marker}.

$readDesc$ is the \texttt{read\_messages} tool description; it gains a
\texttt{(N unread)} marker when mail is pending.  Its source of truth is
$unreadCount$, the server-side count of unread messages, which
\emph{decreases} as mail is read---so a description tied to it clears
itself.  This is the reference channel: it is known to behave correctly,
and it is modelled here so that the talk channel can be judged against a
worked example rather than against prose.

$talkDesc$ is the \texttt{talk} tool description; it gains a
\texttt{[TALK] \ldots} marker, including an accept hint, when a talk
invite, queued message, or live connection exists.  Its activity has
three sources: $talkPending$, a partial function mapping each user who
has invited us to that invite's originating $SESSIONKEY$; $talkQueued$,
the count of undrained talk messages; and $talkConnected$, whether a
conversation is live.  A pending invite in $talkPending$ retains the
inviter's session key, so its accept hint can name a session.

$talkBare$ is the pathological representation: a set of users whose
pending invite was recorded \emph{without} a session key---the tty was
dropped.  A user in $talkBare$ is genuine activity (the marker appears),
but its accept hint can only read \texttt{talk @user} with no
\texttt{:tty}, which is not a runnable accept command.  The structural
invariant $\dom talkPending \cap talkBare = \emptyset$ says a single
invite is recorded one way or the other, never both.  In a correct
implementation $talkBare$ is empty in every reachable state; the model
admits it precisely so the defect can be reached and then proved
unreachable.

\paragraph{The pending-invite clock.}  $pendingAge$ is an abstract clock
for the pending set, bounded by $maxInviteAge$.  It advances on poller
ticks ($AgeTick$) while any invite is pending and is pinned to zero
whenever no invite is pending---the invariant
$(\dom talkPending = \emptyset \land talkBare = \emptyset) \implies
pendingAge = 0$ says exactly this, so a fresh invite always starts its
life at age zero.  Its sole purpose is to make the time-to-live path of
Section~\ref{sec:markerfix} reachable: an invite may be expired
($ExpirePendingInvite$) only once the clock has reached its bound.  The
clock is deliberately global rather than per-invite; at these bounds that
loses nothing the marker property depends on, and it keeps the state
factor at two rather than one per user.

%%----------------------------------------------------------------------
\section{Initialization}
%%----------------------------------------------------------------------

\begin{schema}{Init}
  NotificationState'
  \where
  natsConn' = natsDisconnected \\
  kvProvisioned' = zfalse \\
  talkSubscribed' = zfalse \\
  session' = sessNone \\
  toolContext' = zfalse \\
  wallDesc' = wallDescDefault \\
  talkDesc' = talkDescDefault \\
  readDesc' = readDescDefault \\
  wallDescChanged' = zfalse \\
  talkDescChanged' = zfalse \\
  unreadCount' = 0 \\
  talkPending' = \emptyset \\
  talkBare' = \emptyset \\
  talkQueued' = 0 \\
  talkConnected' = zfalse \\
  pendingAge' = 0 \\
  queueItems' = \langle \rangle \\
  queueKinds' = \emptyset \\
  queueTexts' = \emptyset \\
  currentIndex' = 0 \\
  fileExists' = zfalse \\
  fileUnreadCount' = 0 \\
  fileItems' = \langle \rangle \\
  barLine2' = idleMarker \\
  barStale' = ztrue \\
  lastNotifyResult' = notifyDropped \\
  staleWithdrawHit' = zfalse \\
  pendingNotify' = zfalse \\
  notifyLost' = zfalse \\
  inFlight' = \emptyset \\
  beltPnSnapshot' = zfalse \\
  capPnSnapshot' = zfalse \\
  raceLost' = zfalse
\end{schema}

The initial state has $session' = sessNone$ and
$natsConn' = natsDisconnected$: nothing has connected and nothing has
been notified.

%%----------------------------------------------------------------------
\section{Session Capture}
\label{sec:sessioncapture}
%%----------------------------------------------------------------------

The MCP server lifespan startup captures the session reference
eagerly, before the NATS connection is established and before the
background poller begins.  This is the key operation that prevents
the biff-8g0 bug class: by the time any NATS callback can fire,
the session is guaranteed to have been captured at least once.

$CaptureSession$ below is the case in which there is nothing to
flush---no description changed while the server was starting up---so
the whole operation is the synchronous assignment
\texttt{\_session = session} with no \texttt{await} in it, genuinely
atomic.  When a $pendingNotify$ \emph{is} recorded, capture must also
flush it, and that flush does \texttt{await
session.send\_tool\_list\_changed()}: a real network call, with a real
failure mode and a real gap between deciding to flush and the flush
landing.  Section~\ref{sec:sessionloss} gives that case---shared with
$RecaptureSession$ below, since both call the same
\texttt{capture\_session()}---as $CaptureSessionFlushBegin$,
$CaptureSessionFlushOk$, and $CaptureSessionFlushFail$.

\begin{schema}{CaptureSession}
  \Delta NotificationState
  \where
  pendingAge' = pendingAge \\
  session = sessNone \\
  natsConn = natsDisconnected \\
  pendingNotify = zfalse \\
  session' = sessCaptured \\
  natsConn' = natsConn \\
  lastNotifyResult' = lastNotifyResult \\
  wallDescChanged' = wallDescChanged \\
  talkDescChanged' = talkDescChanged \\
  pendingNotify' = pendingNotify \\
  notifyLost' = notifyLost \\
  kvProvisioned' = kvProvisioned \\
  talkSubscribed' = talkSubscribed \\
  toolContext' = toolContext \\
  readDesc' = readDesc \\
  unreadCount' = unreadCount \\
  talkPending' = talkPending \\
  talkBare' = talkBare \\
  talkQueued' = talkQueued \\
  talkConnected' = talkConnected \\
  wallDesc' = wallDesc \\
  talkDesc' = talkDesc \\
  queueItems' = queueItems \\
  queueKinds' = queueKinds \\
  queueTexts' = queueTexts \\
  currentIndex' = currentIndex \\
  fileExists' = fileExists \\
  fileUnreadCount' = fileUnreadCount \\
  fileItems' = fileItems \\
  barLine2' = barLine2 \\
  barStale' = barStale \\
  inFlight' = inFlight \\
  beltPnSnapshot' = beltPnSnapshot \\
  capPnSnapshot' = capPnSnapshot \\
  raceLost' = raceLost \\
  staleWithdrawHit' = staleWithdrawHit
\end{schema}

After $CaptureSession$, the system has $session = sessCaptured$
with $natsConn = natsDisconnected$, and $EnsureConnected$ can now be
applied.

%%----------------------------------------------------------------------
\section{NATS Connection Operations}
%%----------------------------------------------------------------------

\subsection{Ensure Connected}
\label{sec:ensureconnected}

Lazy connection: transitions from disconnected to connected,
provisioning KV and stream infrastructure.  Idempotent when
already connected.

The guard $session \neq sessNone$ is the biff-8g0 protection, stated
here explicitly rather than derived from a state invariant (see the
erratum in Section~\ref{sec:state}): the system must execute
$CaptureSession$ at least once before the first $EnsureConnected$.
It is deliberately \emph{not} $session = sessCaptured$---after a
suspenders send failure (Section~\ref{sec:sessionloss}) $session$ may
be $sessNone$ again while $natsConn$ is already $natsConnected$, and
$EnsureConnected$ must remain applicable (idempotently) in that state
without requiring a fresh capture first.

\begin{schema}{EnsureConnected}
  \Delta NotificationState
  \where
  pendingAge' = pendingAge \\
  session \neq sessNone \lor natsConn = natsConnected \\
  natsConn = natsDisconnected \lor natsConn = natsConnected \\
  natsConn' = natsConnected \\
  kvProvisioned' = ztrue \\
  talkSubscribed' = talkSubscribed \\
  session' = session \\
  toolContext' = toolContext \\
  readDesc' = readDesc \\
  unreadCount' = unreadCount \\
  talkPending' = talkPending \\
  talkBare' = talkBare \\
  talkQueued' = talkQueued \\
  talkConnected' = talkConnected \\
  wallDesc' = wallDesc \\
  talkDesc' = talkDesc \\
  wallDescChanged' = wallDescChanged \\
  talkDescChanged' = talkDescChanged \\
  queueItems' = queueItems \\
  queueKinds' = queueKinds \\
  queueTexts' = queueTexts \\
  currentIndex' = currentIndex \\
  fileExists' = fileExists \\
  fileUnreadCount' = fileUnreadCount \\
  fileItems' = fileItems \\
  barLine2' = barLine2 \\
  barStale' = barStale \\
  lastNotifyResult' = lastNotifyResult \\
  pendingNotify' = pendingNotify \\
  notifyLost' = notifyLost \\
  inFlight' = inFlight \\
  beltPnSnapshot' = beltPnSnapshot \\
  capPnSnapshot' = capPnSnapshot \\
  raceLost' = raceLost \\
  staleWithdrawHit' = staleWithdrawHit
\end{schema}

\subsection{Close Connection}

Permanent shutdown.  Clears all NATS-dependent state.  The session
reference is preserved by this operation---a clean shutdown does not
by itself invalidate it; only a failed \emph{send} does
(Section~\ref{sec:sessionloss}).

\begin{schema}{CloseConnection}
  \Delta NotificationState
  \where
  pendingAge' = pendingAge \\
  natsConn' = natsClosed \\
  kvProvisioned' = zfalse \\
  talkSubscribed' = zfalse \\
  session' = session \\
  toolContext' = toolContext \\
  readDesc' = readDesc \\
  unreadCount' = unreadCount \\
  talkPending' = talkPending \\
  talkBare' = talkBare \\
  talkQueued' = talkQueued \\
  talkConnected' = talkConnected \\
  wallDesc' = wallDesc \\
  talkDesc' = talkDesc \\
  wallDescChanged' = wallDescChanged \\
  talkDescChanged' = talkDescChanged \\
  queueItems' = queueItems \\
  queueKinds' = queueKinds \\
  queueTexts' = queueTexts \\
  currentIndex' = currentIndex \\
  fileExists' = fileExists \\
  fileUnreadCount' = fileUnreadCount \\
  fileItems' = fileItems \\
  barLine2' = barLine2 \\
  barStale' = barStale \\
  lastNotifyResult' = lastNotifyResult \\
  pendingNotify' = pendingNotify \\
  notifyLost' = notifyLost \\
  inFlight' = inFlight \\
  beltPnSnapshot' = beltPnSnapshot \\
  capPnSnapshot' = capPnSnapshot \\
  raceLost' = raceLost \\
  staleWithdrawHit' = staleWithdrawHit
\end{schema}

%%----------------------------------------------------------------------
\section{Tool Context}
\label{sec:toolcontext}
%%----------------------------------------------------------------------

A tool handler runs \emph{inside} the MCP request-response cycle; the
background poller and the NATS subscription callbacks run \emph{outside}
it.  The component $toolContext$ records which of the two the system is
currently in: $ztrue$ during a tool call, $zfalse$ otherwise.  The
distinction is not decoration.  It selects the notification path---belt
when a request context is available, suspenders when it is not---and it
gates every operation that can only occur inside a handler ($WallPost$,
$TalkStart$, $NotifyBelt$) against every operation that can only occur
in the background ($NotifySuspenders$, $KVWallReceive$,
$NatsTalkCallback$, $MessagePushCallback$, $PollTick$).

The earlier drafts left $toolContext$ at its initial $zfalse$ with no
operation to raise it.  The consequence was quiet but total: no reachable
state had $toolContext = ztrue$, so the belt-path operations were enabled
in no reachable state, and the guarantees stated for them
($NotifyGuarantee$, $WallNotifySuccess$) held only \emph{vacuously}.  A
vacuous proof is not a proof of anything.  The remedy is to model the
entry into and exit from a tool handler explicitly, so that the belt path
is reachable and its obligations bite.

$EnterToolContext$ marks the start of a tool call; $LeaveToolContext$
marks its end.  Each toggles $toolContext$ and changes nothing else, so
neither can disturb any invariant.  Their guards make the two mutually
exclusive: one may enter only from the background, and leave only from a
handler.

\begin{schema}{EnterToolContext}
  \Delta NotificationState
  \where
  pendingAge' = pendingAge \\
  toolContext = zfalse \\
  toolContext' = ztrue \\
  natsConn' = natsConn \\
  kvProvisioned' = kvProvisioned \\
  talkSubscribed' = talkSubscribed \\
  session' = session \\
  wallDesc' = wallDesc \\
  talkDesc' = talkDesc \\
  readDesc' = readDesc \\
  wallDescChanged' = wallDescChanged \\
  talkDescChanged' = talkDescChanged \\
  unreadCount' = unreadCount \\
  talkPending' = talkPending \\
  talkBare' = talkBare \\
  talkQueued' = talkQueued \\
  talkConnected' = talkConnected \\
  queueItems' = queueItems \\
  queueKinds' = queueKinds \\
  queueTexts' = queueTexts \\
  currentIndex' = currentIndex \\
  fileExists' = fileExists \\
  fileUnreadCount' = fileUnreadCount \\
  fileItems' = fileItems \\
  barLine2' = barLine2 \\
  barStale' = barStale \\
  lastNotifyResult' = lastNotifyResult \\
  pendingNotify' = pendingNotify \\
  notifyLost' = notifyLost \\
  inFlight' = inFlight \\
  beltPnSnapshot' = beltPnSnapshot \\
  capPnSnapshot' = capPnSnapshot \\
  raceLost' = raceLost \\
  staleWithdrawHit' = staleWithdrawHit
\end{schema}

\begin{schema}{LeaveToolContext}
  \Delta NotificationState
  \where
  pendingAge' = pendingAge \\
  toolContext = ztrue \\
  toolContext' = zfalse \\
  natsConn' = natsConn \\
  kvProvisioned' = kvProvisioned \\
  talkSubscribed' = talkSubscribed \\
  session' = session \\
  wallDesc' = wallDesc \\
  talkDesc' = talkDesc \\
  readDesc' = readDesc \\
  wallDescChanged' = wallDescChanged \\
  talkDescChanged' = talkDescChanged \\
  unreadCount' = unreadCount \\
  talkPending' = talkPending \\
  talkBare' = talkBare \\
  talkQueued' = talkQueued \\
  talkConnected' = talkConnected \\
  queueItems' = queueItems \\
  queueKinds' = queueKinds \\
  queueTexts' = queueTexts \\
  currentIndex' = currentIndex \\
  fileExists' = fileExists \\
  fileUnreadCount' = fileUnreadCount \\
  fileItems' = fileItems \\
  barLine2' = barLine2 \\
  barStale' = barStale \\
  lastNotifyResult' = lastNotifyResult \\
  pendingNotify' = pendingNotify \\
  notifyLost' = notifyLost \\
  inFlight' = inFlight \\
  beltPnSnapshot' = beltPnSnapshot \\
  capPnSnapshot' = capPnSnapshot \\
  raceLost' = raceLost \\
  staleWithdrawHit' = staleWithdrawHit
\end{schema}

A tool handler is therefore modelled as $EnterToolContext$, followed by
the handler's own operation on the belt path, followed by
$LeaveToolContext$; background work happens between these brackets, never
inside them.  This is the smallest change that makes the belt path
reachable without weakening any existing guard.

%%----------------------------------------------------------------------
\section{Notification Core}
\label{sec:notifycore}
%%----------------------------------------------------------------------

This is the central mechanism.  Two paths send the MCP
\texttt{tools/list\_changed} notification.  \emph{Erratum}: earlier
drafts stated the belt path ``always succeeds.''  Round~2's review
found it does not---$NotifyBeltSendFail$
(Section~\ref{sec:beltpath})---so both paths can fail, and both are
given here in the simplified, idealised form they were first
understood in, before Sections~\ref{sec:sessionloss}
and~\ref{sec:capflush} model the failure at each site precisely and
Section~\ref{sec:race} models the three sites' non-atomicity.

\subsection{Belt Path (Tool Context Available)}
\label{sec:beltpath}

Inside a tool handler: notification piggybacks on the MCP response.
This is the preferred path when tool context is available because it
guarantees synchronous delivery within the request-response cycle---but
``synchronous'' describes the request-response cycle, not the call
itself: \texttt{await ctx.send\_notification(...)} is still a genuine
\texttt{await}, and round~2's review found it can fail with something
other than the \texttt{RuntimeError} that selects belt over suspenders
in the first place---a transport error from a client that dropped mid-request.
The real implementation also refreshes the stored session reference on
\emph{every} belt call, whether or not the prior reference was still
good---\texttt{\_session = ctx.session}, unconditionally, and only once
the send has actually succeeded---so a belt call is a critical section
with the same read/decide-then-await-then-write shape as the other two
sites.  $NotifyBeltBegin$ is the read/decide step; $NotifyBelt$ and
$NotifyBeltSendFail$ are its two possible commits.

\begin{schema}{NotifyBeltBegin}
  \Delta NotificationState
  \where
  pendingAge' = pendingAge \\
  toolContext = ztrue \\
  siteBelt \notin inFlight \\
  inFlight' = inFlight \cup \{ siteBelt \} \\
  beltPnSnapshot' = pendingNotify \\
  session' = session \\
  natsConn' = natsConn \\
  kvProvisioned' = kvProvisioned \\
  talkSubscribed' = talkSubscribed \\
  toolContext' = toolContext \\
  readDesc' = readDesc \\
  unreadCount' = unreadCount \\
  talkPending' = talkPending \\
  talkBare' = talkBare \\
  talkQueued' = talkQueued \\
  talkConnected' = talkConnected \\
  wallDesc' = wallDesc \\
  talkDesc' = talkDesc \\
  wallDescChanged' = wallDescChanged \\
  talkDescChanged' = talkDescChanged \\
  queueItems' = queueItems \\
  queueKinds' = queueKinds \\
  queueTexts' = queueTexts \\
  currentIndex' = currentIndex \\
  fileExists' = fileExists \\
  fileUnreadCount' = fileUnreadCount \\
  fileItems' = fileItems \\
  barLine2' = barLine2 \\
  barStale' = barStale \\
  lastNotifyResult' = lastNotifyResult \\
  pendingNotify' = pendingNotify \\
  notifyLost' = notifyLost \\
  capPnSnapshot' = capPnSnapshot \\
  raceLost' = raceLost \\
  staleWithdrawHit' = staleWithdrawHit
\end{schema}

$NotifyBelt$ is the send-succeeds commit.  It always flushes---clearing
$pendingNotify$ unconditionally is correct precisely because the belt
call's own send just delivered the \emph{current} tool description,
which reflects every mutation up to this point, whichever site armed
$pendingNotify$.  $raceLost$ is Section~\ref{sec:race}'s instrument: if
$pendingNotify$ has changed since this call's own $beltPnSnapshot$ was
taken at $NotifyBeltBegin$, some other site's write landed in the gap,
and this commit is about to overwrite it with a decision made before
that write happened.

\begin{schema}{NotifyBelt}
  \Delta NotificationState
  \where
  pendingAge' = pendingAge \\
  toolContext = ztrue \\
  siteBelt \in inFlight \\
  inFlight' = inFlight \setminus \{ siteBelt \} \\
  lastNotifyResult' = notifySent \\
  session' = sessCaptured \\
  pendingNotify' = zfalse \\
  ((pendingNotify \neq beltPnSnapshot) \implies raceLost' = ztrue) \\
  ((pendingNotify = beltPnSnapshot) \implies raceLost' = raceLost) \\
  natsConn' = natsConn \\
  kvProvisioned' = kvProvisioned \\
  talkSubscribed' = talkSubscribed \\
  toolContext' = toolContext \\
  readDesc' = readDesc \\
  unreadCount' = unreadCount \\
  talkPending' = talkPending \\
  talkBare' = talkBare \\
  talkQueued' = talkQueued \\
  talkConnected' = talkConnected \\
  wallDesc' = wallDesc \\
  talkDesc' = talkDesc \\
  wallDescChanged' = zfalse \\
  talkDescChanged' = zfalse \\
  queueItems' = queueItems \\
  queueKinds' = queueKinds \\
  queueTexts' = queueTexts \\
  currentIndex' = currentIndex \\
  fileExists' = fileExists \\
  fileUnreadCount' = fileUnreadCount \\
  fileItems' = fileItems \\
  barLine2' = barLine2 \\
  barStale' = zfalse \\
  notifyLost' = notifyLost \\
  beltPnSnapshot' = beltPnSnapshot \\
  capPnSnapshot' = capPnSnapshot \\
  staleWithdrawHit' = staleWithdrawHit
\end{schema}

$NotifyBeltSendFail$ is the corrected failure commit: the send raised
something other than \texttt{RuntimeError}, nothing was delivered, the
session assignment line was never reached (it comes \emph{after} the
\texttt{await} in the real code, so $session$ is left exactly as it
was), and---the round~2 fix---$pendingNotify$ is raised to record that
this attempt failed.  $NotifyBeltSendFailUnrecorded$ is the dead twin:
identical except $pendingNotify$ is left untouched, which is exactly
what happens if the exception handler that round~2 asks for is never
written.

\begin{schema}{NotifyBeltSendFail}
  \Delta NotificationState
  \where
  pendingAge' = pendingAge \\
  toolContext = ztrue \\
  siteBelt \in inFlight \\
  inFlight' = inFlight \setminus \{ siteBelt \} \\
  lastNotifyResult' = notifyDropped \\
  session' = session \\
  pendingNotify' = ztrue \\
  natsConn' = natsConn \\
  kvProvisioned' = kvProvisioned \\
  talkSubscribed' = talkSubscribed \\
  toolContext' = toolContext \\
  readDesc' = readDesc \\
  unreadCount' = unreadCount \\
  talkPending' = talkPending \\
  talkBare' = talkBare \\
  talkQueued' = talkQueued \\
  talkConnected' = talkConnected \\
  wallDesc' = wallDesc \\
  talkDesc' = talkDesc \\
  wallDescChanged' = wallDescChanged \\
  talkDescChanged' = talkDescChanged \\
  queueItems' = queueItems \\
  queueKinds' = queueKinds \\
  queueTexts' = queueTexts \\
  currentIndex' = currentIndex \\
  fileExists' = fileExists \\
  fileUnreadCount' = fileUnreadCount \\
  fileItems' = fileItems \\
  barLine2' = barLine2 \\
  barStale' = barStale \\
  notifyLost' = notifyLost \\
  beltPnSnapshot' = beltPnSnapshot \\
  capPnSnapshot' = capPnSnapshot \\
  raceLost' = raceLost \\
  staleWithdrawHit' = staleWithdrawHit
\end{schema}

\begin{schema}{NotifyBeltSendFailUnrecorded}
  \Delta NotificationState
  \where
  pendingAge' = pendingAge \\
  toolContext = ztrue \\
  siteBelt \in inFlight \\
  inFlight' = inFlight \setminus \{ siteBelt \} \\
  lastNotifyResult' = notifyDropped \\
  session' = session \\
  pendingNotify' = pendingNotify \\
  notifyLost' = ztrue \\
  natsConn' = natsConn \\
  kvProvisioned' = kvProvisioned \\
  talkSubscribed' = talkSubscribed \\
  toolContext' = toolContext \\
  readDesc' = readDesc \\
  unreadCount' = unreadCount \\
  talkPending' = talkPending \\
  talkBare' = talkBare \\
  talkQueued' = talkQueued \\
  talkConnected' = talkConnected \\
  wallDesc' = wallDesc \\
  talkDesc' = talkDesc \\
  wallDescChanged' = wallDescChanged \\
  talkDescChanged' = talkDescChanged \\
  queueItems' = queueItems \\
  queueKinds' = queueKinds \\
  queueTexts' = queueTexts \\
  currentIndex' = currentIndex \\
  fileExists' = fileExists \\
  fileUnreadCount' = fileUnreadCount \\
  fileItems' = fileItems \\
  barLine2' = barLine2 \\
  barStale' = barStale \\
  beltPnSnapshot' = beltPnSnapshot \\
  capPnSnapshot' = capPnSnapshot \\
  raceLost' = raceLost \\
  staleWithdrawHit' = staleWithdrawHit
\end{schema}

\subsection{Suspenders Path (No Tool Context), idealised}

Background poller: no request context, so the eagerly captured
session reference is used to send the notification instead.  NATS
callbacks do not use this path directly---they set description-changed
flags that the poller detects on its next tick.  This schema is the
\emph{idealised} suspenders path, kept for continuity with the earlier
drafts: it assumes the send always succeeds whenever
$session = sessCaptured$.  Section~\ref{sec:sessionloss} replaces this
assumption at the operational level with the two outcomes a real send
actually has.

\begin{schema}{NotifySuspenders}
  \Delta NotificationState
  \where
  pendingAge' = pendingAge \\
  toolContext = zfalse \\
  session = sessCaptured \\
  lastNotifyResult' = notifySent \\
  session' = session \\
  natsConn' = natsConn \\
  kvProvisioned' = kvProvisioned \\
  talkSubscribed' = talkSubscribed \\
  toolContext' = toolContext \\
  readDesc' = readDesc \\
  unreadCount' = unreadCount \\
  talkPending' = talkPending \\
  talkBare' = talkBare \\
  talkQueued' = talkQueued \\
  talkConnected' = talkConnected \\
  wallDesc' = wallDesc \\
  talkDesc' = talkDesc \\
  wallDescChanged' = zfalse \\
  talkDescChanged' = zfalse \\
  queueItems' = queueItems \\
  queueKinds' = queueKinds \\
  queueTexts' = queueTexts \\
  currentIndex' = currentIndex \\
  fileExists' = fileExists \\
  fileUnreadCount' = fileUnreadCount \\
  fileItems' = fileItems \\
  barLine2' = barLine2 \\
  barStale' = zfalse \\
  pendingNotify' = pendingNotify \\
  notifyLost' = notifyLost \\
  inFlight' = inFlight \\
  beltPnSnapshot' = beltPnSnapshot \\
  capPnSnapshot' = capPnSnapshot \\
  raceLost' = raceLost \\
  staleWithdrawHit' = staleWithdrawHit
\end{schema}

\subsection{Total Notification}

The complete notification operation is the disjunction of the two
paths.  Exactly one fires: $NotifyBelt$ when $toolContext = ztrue$,
$NotifySuspenders$ (idealised) or its two-outcome refinement,
Section~\ref{sec:sessionloss}, when $toolContext = zfalse$.

\paragraph{Erratum.}  The paragraph that stood here in earlier drafts
claimed ``no notification is ever dropped,'' reasoning from the state
invariant $natsConn = natsConnected \implies session = sessCaptured$
that Section~\ref{sec:state}'s erratum has since retracted.  That
claim is false: a suspenders send can fail independently of NATS
connectivity, and root cause~\#1 (Section~\ref{sec:sessionloss}) is
exactly the case in which the resulting drop was, until the fix given
there, never recorded and never recovered.  NATS callbacks still defer
to the poller and wake it from napping so that the next active tick
\emph{attempts} the notification promptly; whether that attempt
succeeds is the question Section~\ref{sec:sessionloss} answers.

%%----------------------------------------------------------------------
\section{Send Flow Operations}
%%----------------------------------------------------------------------

\subsection{Wall Post (Tool Handler)}

The wall tool executes inside a tool handler context.  It
mutates the wall description, adds a display item, and fires
notification via the belt path.  The precondition requires
tool context and a live NATS connection.

\begin{schema}{WallPost}
  \Delta NotificationState \\
  wallText? : MSG \\
  wallKey? : SOURCEKEY
  \where
  pendingAge' = pendingAge \\
  toolContext = ztrue \\
  natsConn = natsConnected \\
  wallDesc' \neq wallDesc \\
  wallDescChanged' = ztrue \\
  (wallKey? \notin \ran queueItems \implies \\
  \quad~ queueItems' = queueItems \cat \langle wallKey? \rangle \land \\
  \quad~ queueKinds' = queueKinds \cup \\
  \quad~ \quad~ \{ wallKey? \mapsto itemWall \} \land \\
  \quad~ queueTexts' = queueTexts \cup \\
  \quad~ \quad~ \{ wallKey? \mapsto wallText? \}) \\
  (wallKey? \in \ran queueItems \implies \\
  \quad~ queueItems' = queueItems \land \\
  \quad~ queueKinds' = queueKinds \land \\
  \quad~ queueTexts' = queueTexts) \\
  session' = session \\
  lastNotifyResult' = notifySent \\
  barStale' = zfalse \\
  natsConn' = natsConn \\
  kvProvisioned' = kvProvisioned \\
  talkSubscribed' = talkSubscribed \\
  toolContext' = toolContext \\
  readDesc' = readDesc \\
  unreadCount' = unreadCount \\
  talkPending' = talkPending \\
  talkBare' = talkBare \\
  talkQueued' = talkQueued \\
  talkConnected' = talkConnected \\
  talkDesc' = talkDesc \\
  talkDescChanged' = zfalse \\
  currentIndex' = currentIndex \\
  fileExists' = fileExists \\
  fileUnreadCount' = fileUnreadCount \\
  fileItems' = fileItems \\
  barLine2' = barLine2 \\
  pendingNotify' = pendingNotify \\
  notifyLost' = notifyLost \\
  inFlight' = inFlight \\
  beltPnSnapshot' = beltPnSnapshot \\
  capPnSnapshot' = capPnSnapshot \\
  raceLost' = raceLost \\
  staleWithdrawHit' = staleWithdrawHit
\end{schema}

\subsection{Talk Start (Tool Handler)}

The talk tool initiates a conversation.  Runs inside tool context,
and arranges for the background poller to subscribe to NATS talk
notifications.

\begin{schema}{TalkStart}
  \Delta NotificationState \\
  partner? : USER
  \where
  pendingAge' = pendingAge \\
  toolContext = ztrue \\
  natsConn = natsConnected \\
  session' = session \\
  lastNotifyResult' = notifySent \\
  talkSubscribed' = ztrue \\
  barStale' = zfalse \\
  natsConn' = natsConn \\
  kvProvisioned' = kvProvisioned \\
  toolContext' = toolContext \\
  readDesc' = readDesc \\
  unreadCount' = unreadCount \\
  talkPending' = talkPending \\
  talkBare' = talkBare \\
  talkQueued' = talkQueued \\
  talkConnected' = talkConnected \\
  wallDesc' = wallDesc \\
  wallDescChanged' = zfalse \\
  talkDesc' = talkDesc \\
  talkDescChanged' = zfalse \\
  queueItems' = queueItems \\
  queueKinds' = queueKinds \\
  queueTexts' = queueTexts \\
  currentIndex' = currentIndex \\
  fileExists' = fileExists \\
  fileUnreadCount' = fileUnreadCount \\
  fileItems' = fileItems \\
  barLine2' = barLine2 \\
  pendingNotify' = pendingNotify \\
  notifyLost' = notifyLost \\
  inFlight' = inFlight \\
  beltPnSnapshot' = beltPnSnapshot \\
  capPnSnapshot' = capPnSnapshot \\
  raceLost' = raceLost \\
  staleWithdrawHit' = staleWithdrawHit
\end{schema}

\subsection{External Message Deliver (CLI)}

A CLI client or another MCP server delivers a message via NATS.
This publishes to JetStream (reliable) and then publishes a
talk notification on core NATS (fire-and-forget).  No tool
context exists on the sender side.  The receiver's $NatsTalkCallback$
handles the receive-side effects.

\begin{schema}{ExternalDeliver}
  \Delta NotificationState \\
  body? : MSG \\
  toUser? : USER
  \where
  pendingAge' = pendingAge \\
  natsConn = natsConnected \\
  natsConn' = natsConn \\
  kvProvisioned' = kvProvisioned \\
  talkSubscribed' = talkSubscribed \\
  session' = session \\
  toolContext' = toolContext \\
  readDesc' = readDesc \\
  unreadCount' = unreadCount \\
  talkPending' = talkPending \\
  talkBare' = talkBare \\
  talkQueued' = talkQueued \\
  talkConnected' = talkConnected \\
  wallDesc' = wallDesc \\
  talkDesc' = talkDesc \\
  wallDescChanged' = wallDescChanged \\
  talkDescChanged' = talkDescChanged \\
  queueItems' = queueItems \\
  queueKinds' = queueKinds \\
  queueTexts' = queueTexts \\
  currentIndex' = currentIndex \\
  fileExists' = fileExists \\
  fileUnreadCount' = fileUnreadCount \\
  fileItems' = fileItems \\
  barLine2' = barLine2 \\
  barStale' = barStale \\
  lastNotifyResult' = lastNotifyResult \\
  pendingNotify' = pendingNotify \\
  notifyLost' = notifyLost \\
  inFlight' = inFlight \\
  beltPnSnapshot' = beltPnSnapshot \\
  capPnSnapshot' = capPnSnapshot \\
  raceLost' = raceLost \\
  staleWithdrawHit' = staleWithdrawHit
\end{schema}

%%----------------------------------------------------------------------
\section{Receive Flow Operations}
%%----------------------------------------------------------------------

\subsection{KV Wall Receive}

Models the effects that occur when the periodic status poller, running
in the MCP server coroutine context, detects that the wall KV entry
has changed (typically after being woken by the NATS KV watcher).
The watcher itself runs \emph{outside} any tool handler context in a
background asyncio task and only wakes the poller; it does \emph{not}
update the wall description, modify the display queue, or fire an MCP
notification directly---sending notifications from a background task
is unreliable for the same reason as NATS subscription callbacks.

After being woken, the poller performs its next active tick (every
2~seconds), reads the wall value from KV, detects any change, applies
this schema to update the wall description and display queue, and then
fires the notification from the MCP server's coroutine context.

The precondition $natsConn = natsConnected$ implies
$session = sessCaptured$ (by state invariant), so the poller will
have a valid session when it fires the deferred notification.

\begin{schema}{KVWallReceive}
  \Delta NotificationState \\
  wallText? : MSG \\
  wallKey? : SOURCEKEY
  \where
  pendingAge' = pendingAge \\
  natsConn = natsConnected \\
  toolContext = zfalse \\
  wallDesc' \neq wallDesc \\
  wallDescChanged' = ztrue \\
  (wallKey? \notin \ran queueItems \implies \\
  \quad~ queueItems' = queueItems \cat \langle wallKey? \rangle \land \\
  \quad~ queueKinds' = queueKinds \cup \\
  \quad~ \quad~ \{ wallKey? \mapsto itemWall \} \land \\
  \quad~ queueTexts' = queueTexts \cup \\
  \quad~ \quad~ \{ wallKey? \mapsto wallText? \}) \\
  (wallKey? \in \ran queueItems \implies \\
  \quad~ queueItems' = queueItems \land \\
  \quad~ queueKinds' = queueKinds \land \\
  \quad~ queueTexts' = queueTexts) \\
  lastNotifyResult' = lastNotifyResult \\
  barStale' = barStale \\
  session' = session \\
  natsConn' = natsConn \\
  kvProvisioned' = kvProvisioned \\
  talkSubscribed' = talkSubscribed \\
  toolContext' = toolContext \\
  readDesc' = readDesc \\
  unreadCount' = unreadCount \\
  talkPending' = talkPending \\
  talkBare' = talkBare \\
  talkQueued' = talkQueued \\
  talkConnected' = talkConnected \\
  talkDesc' = talkDesc \\
  talkDescChanged' = talkDescChanged \\
  currentIndex' = currentIndex \\
  fileExists' = fileExists \\
  fileUnreadCount' = fileUnreadCount \\
  fileItems' = fileItems \\
  barLine2' = barLine2 \\
  pendingNotify' = pendingNotify \\
  notifyLost' = notifyLost \\
  inFlight' = inFlight \\
  beltPnSnapshot' = beltPnSnapshot \\
  capPnSnapshot' = capPnSnapshot \\
  raceLost' = raceLost \\
  staleWithdrawHit' = staleWithdrawHit
\end{schema}

\subsection{NATS Talk Callback}

Fires when a core NATS message arrives on the talk notification
subject.  Runs \emph{outside} any tool handler context in a NATS
subscription callback coroutine.  Mutates the talk description and
adds the talk item to the display queue, but does \emph{not} fire
an MCP notification directly---sending notifications from a NATS
callback is unreliable because Claude~Code intermittently ignores
them (different coroutine context from the MCP server).

Instead, the callback sets $talkDescChanged' = ztrue$ and the
background poller ($PollTick$) detects the change on its next tick
and fires the notification from the MCP server's own coroutine
context.  This is the same pattern wall uses: state mutation is
decoupled from notification delivery.

The precondition $talkSubscribed = ztrue$ implies
$natsConn = natsConnected$ (by state invariant), which in turn implies
$session = sessCaptured$.  The poller will therefore have a valid
session when it fires the deferred notification.

\begin{schema}{NatsTalkCallback}
  \Delta NotificationState \\
  senderBody? : MSG \\
  senderKey? : SOURCEKEY
  \where
  pendingAge' = pendingAge \\
  talkSubscribed = ztrue \\
  toolContext = zfalse \\
  talkBare = \emptyset \\
  talkQueued < maxQueueSize \\
  talkDesc' \neq talkDescDefault \\
  talkDescChanged' = ztrue \\
  queueItems' = queueItems \cat \langle senderKey? \rangle \\
  queueKinds' = queueKinds \oplus \\
  \quad~ \{ senderKey? \mapsto itemTalk \} \\
  queueTexts' = queueTexts \oplus \\
  \quad~ \{ senderKey? \mapsto senderBody? \} \\
  currentIndex' = \# queueItems \\
  lastNotifyResult' = lastNotifyResult \\
  barStale' = barStale \\
  session' = session \\
  natsConn' = natsConn \\
  kvProvisioned' = kvProvisioned \\
  talkSubscribed' = talkSubscribed \\
  toolContext' = toolContext \\
  readDesc' = readDesc \\
  unreadCount' = unreadCount \\
  talkPending' = talkPending \\
  talkBare' = talkBare \\
  talkQueued' = talkQueued + 1 \\
  talkConnected' = talkConnected \\
  wallDesc' = wallDesc \\
  wallDescChanged' = wallDescChanged \\
  fileExists' = fileExists \\
  fileUnreadCount' = fileUnreadCount \\
  fileItems' = fileItems \\
  barLine2' = barLine2 \\
  pendingNotify' = pendingNotify \\
  notifyLost' = notifyLost \\
  inFlight' = inFlight \\
  beltPnSnapshot' = beltPnSnapshot \\
  capPnSnapshot' = capPnSnapshot \\
  raceLost' = raceLost \\
  staleWithdrawHit' = staleWithdrawHit
\end{schema}

\subsection{Message Push Callback}

Fires when a core NATS message arrives on the broadcast inbox-notify
subject (DES-062, \texttt{nats-relay.tex}'s $inboxNotify$ kind of $subGen$).
Runs \emph{outside} any tool handler context, exactly like
$NatsTalkCallback$ and $KVWallReceive$---but unlike those two, the poke
carries no payload: it is the bare fact that \emph{something} changed on
the durable JetStream inbox, not a copy of what changed.  A broadcast
message has no session identity to target (talk.tex $subjectOf$ requires
a $user:tty$ recipient), so the wake poke names only the bare repo+user
pair the durable inbox is partitioned by---there is nothing in the frame
for this callback to decode, and therefore nothing it can write into
$unreadCount$ or $readDesc$.  Recomputing those still requires the actual
$stream\_info$ round-trip, which only the poller performs, on its next
tick.

\begin{schema}{MessagePushCallback}
  \Xi NotificationState
  \where
  natsConn = natsConnected \\
  toolContext = zfalse
\end{schema}

$MessagePushCallback$ is the identity transition on every modelled
variable---$unreadCount$, $readDesc$, $pendingNotify$, and $notifyLost$
included---because it never calls \texttt{refresh\_read\_messages} or
\texttt{notify\_tool\_list\_changed} directly (DES-020/DES-021); the code
side of this callback does exactly one visible-to-the-model-adjacent
thing, waking the poller out of napping, and that wake lives entirely
outside $NotificationState$'s vocabulary---the poller's active/napping
timer is a separate, untimed mechanism the model does not represent, the
same way $KVWallReceive$'s and $NatsTalkCallback$'s own wake effect is
already out of scope here.  Because $MessagePushCallback$ changes nothing,
adding it falsifies no existing property: $NotifyGuarantee$ and
$notifyLost = zfalse$ (Section~\ref{sec:sessionloss}) both range only over
the states $NatsTalkCallback$/$KVWallReceive$/$PollTick$ already reach, and
a purely additive self-loop operation cannot remove a state, an edge, or a
counterexample-witnessing path any existing $AG$/$EF$ property depends on.
This is the mechanical extension the design anticipated
(\texttt{design-nats-push.md} \S4): a third ``wake, defer the send'' source,
following the $NatsTalkCallback$/$KVWallReceive$ template, added with no
reconciliation of any conflicting model or code because none exists.

\subsection{Sync Unread File}

Writes the display queue snapshot to the filesystem status file.
Called after any state change that affects what the status bar shows.

\begin{schema}{SyncUnreadFile}
  \Delta NotificationState \\
  count? : \nat
  \where
  pendingAge' = pendingAge \\
  count? \leq maxUnreadCount \\
  fileExists' = ztrue \\
  fileUnreadCount' = count? \\
  fileItems' = queueItems \\
  natsConn' = natsConn \\
  kvProvisioned' = kvProvisioned \\
  talkSubscribed' = talkSubscribed \\
  session' = session \\
  toolContext' = toolContext \\
  readDesc' = readDesc \\
  unreadCount' = unreadCount \\
  talkPending' = talkPending \\
  talkBare' = talkBare \\
  talkQueued' = talkQueued \\
  talkConnected' = talkConnected \\
  wallDesc' = wallDesc \\
  talkDesc' = talkDesc \\
  wallDescChanged' = wallDescChanged \\
  talkDescChanged' = talkDescChanged \\
  queueItems' = queueItems \\
  queueKinds' = queueKinds \\
  queueTexts' = queueTexts \\
  currentIndex' = currentIndex \\
  barLine2' = barLine2 \\
  barStale' = barStale \\
  lastNotifyResult' = lastNotifyResult \\
  pendingNotify' = pendingNotify \\
  notifyLost' = notifyLost \\
  inFlight' = inFlight \\
  beltPnSnapshot' = beltPnSnapshot \\
  capPnSnapshot' = capPnSnapshot \\
  raceLost' = raceLost \\
  staleWithdrawHit' = staleWithdrawHit
\end{schema}

\subsection{Poll Tick (Background Poller)}
\label{sec:polltick}

The 2-second background poller tick.  Runs outside tool context.  This
schema is the \emph{idle} case only---nothing has changed, so there is
nothing to attempt.  Section~\ref{sec:sessionloss} gives the three
operations that cover the remaining case, $wallDescChanged = ztrue \lor
talkDescChanged = ztrue$: a successful suspenders send, a failed one
that is correctly recorded, and---kept as the now-dead root cause
\#1---a failed one that is silently absorbed.

\begin{schema}{PollTick}
  \Delta NotificationState
  \where
  pendingAge' = pendingAge \\
  toolContext = zfalse \\
  wallDescChanged = zfalse \\
  talkDescChanged = zfalse \\
  lastNotifyResult' = lastNotifyResult \\
  barStale' = barStale \\
  natsConn' = natsConn \\
  kvProvisioned' = kvProvisioned \\
  talkSubscribed' = talkSubscribed \\
  session' = session \\
  toolContext' = toolContext \\
  readDesc' = readDesc \\
  unreadCount' = unreadCount \\
  talkPending' = talkPending \\
  talkBare' = talkBare \\
  talkQueued' = talkQueued \\
  talkConnected' = talkConnected \\
  wallDesc' = wallDesc \\
  talkDesc' = talkDesc \\
  wallDescChanged' = zfalse \\
  talkDescChanged' = zfalse \\
  queueItems' = queueItems \\
  queueKinds' = queueKinds \\
  queueTexts' = queueTexts \\
  currentIndex' = currentIndex \\
  fileExists' = fileExists \\
  fileUnreadCount' = fileUnreadCount \\
  fileItems' = fileItems \\
  barLine2' = barLine2 \\
  pendingNotify' = pendingNotify \\
  notifyLost' = notifyLost \\
  inFlight' = inFlight \\
  beltPnSnapshot' = beltPnSnapshot \\
  capPnSnapshot' = capPnSnapshot \\
  raceLost' = raceLost \\
  staleWithdrawHit' = staleWithdrawHit
\end{schema}

%%----------------------------------------------------------------------
\section{Suspenders Failure and the No-Recovery Sink}
\label{sec:sessionloss}
%%----------------------------------------------------------------------

Root cause~\#1, ``dropped \texttt{tools/list\_changed} mid-session, no
recovery.''  A suspenders send can fail---the stored
\texttt{ServerSession} may be attached to a stream the client has
already torn down---independently of whether the NATS connection
itself is healthy.  $notify\_tool\_list\_changed()$'s suspenders branch
catches exactly this and clears the stored session reference so the
poller stops hammering a dead stream:

\begin{verbatim}
if _session is not None:
    try:
        await _session.send_tool_list_changed()
    except Exception:
        logger.warning(...)
        _session = None
    return
_pending_notify = True
\end{verbatim}

The \texttt{return} on the exception path is the defect: it skips the
\texttt{\_pending\_notify = True} line entirely.  A failed suspenders
send and a session that was simply never captured end up
indistinguishable---both leave $\texttt{\_session} = \texttt{None}$
with nothing recorded---even though the two cases are not
symmetric.  The pre-session case is genuinely a first attempt with
nothing yet to notify about; the suspenders-failure case has already
mutated the tool description and already decided a notification is
owed.  Losing that distinction is what makes the drop permanent: the
one flag that exists to remember an owed notification, $pendingNotify$,
is never set.

$PollTickNotifyBegin$ is the read/decide step: a changed flag is set,
the session looks captured, so the poller commits to attempting a
send.  Nothing about the send's outcome is known yet.

\begin{schema}{PollTickNotifyBegin}
  \Delta NotificationState
  \where
  pendingAge' = pendingAge \\
  toolContext = zfalse \\
  (wallDescChanged = ztrue \lor talkDescChanged = ztrue) \\
  session = sessCaptured \\
  siteSuspenders \notin inFlight \\
  inFlight' = inFlight \cup \{ siteSuspenders \} \\
  session' = session \\
  natsConn' = natsConn \\
  kvProvisioned' = kvProvisioned \\
  talkSubscribed' = talkSubscribed \\
  toolContext' = toolContext \\
  readDesc' = readDesc \\
  unreadCount' = unreadCount \\
  talkPending' = talkPending \\
  talkBare' = talkBare \\
  talkQueued' = talkQueued \\
  talkConnected' = talkConnected \\
  wallDesc' = wallDesc \\
  talkDesc' = talkDesc \\
  wallDescChanged' = wallDescChanged \\
  talkDescChanged' = talkDescChanged \\
  queueItems' = queueItems \\
  queueKinds' = queueKinds \\
  queueTexts' = queueTexts \\
  currentIndex' = currentIndex \\
  fileExists' = fileExists \\
  fileUnreadCount' = fileUnreadCount \\
  fileItems' = fileItems \\
  barLine2' = barLine2 \\
  barStale' = barStale \\
  lastNotifyResult' = lastNotifyResult \\
  pendingNotify' = pendingNotify \\
  notifyLost' = notifyLost \\
  beltPnSnapshot' = beltPnSnapshot \\
  capPnSnapshot' = capPnSnapshot \\
  raceLost' = raceLost \\
  staleWithdrawHit' = staleWithdrawHit
\end{schema}

$PollTickNotifyOk$ commits a successful send: releases $siteSuspenders$
from $inFlight$, clears the changed flags, and---the biff-6vuv
amendment---also clears $pendingNotify$.  Earlier drafts left
$pendingNotify$ untouched here on the argument that ``the suspenders path
never has anything of its own to flush, only the current tick's own
delivery to report''; that argument proves too little.  $NotifyBelt$
(Section~\ref{sec:beltpath}) clears $pendingNotify$ unconditionally for
exactly the reason restated here: whichever site armed the flag, a send
that just succeeded delivered the \emph{current} tool description, which
reflects every mutation up to this point---so the debt the flag was
recording is paid, regardless of which site incurred it.  $PollTickNotifyOk$
is a successful send by the same measure $NotifyBelt$ is; withholding the
clear only because this send happened to run on the suspenders path
re-armed nothing new and healed nothing---it left a stale $ztrue$ sitting
on a debt this very commit already discharged, so an unrelated future
belt call would eventually flush a notification for a description that had
already gone out.  No snapshot-based race instrumentation
($beltPnSnapshot$/$capPnSnapshot$/$raceLost$) is added for this site: those
exist for $siteBelt$ and $siteCapture$ because their commits are conditioned
on a value read at their own $\dots Begin$ step, and a competing write could
land in the gap between; $PollTickNotifyOk$'s clear is unconditional, reads
nothing back, and races nothing by construction, and the
$\# inFlight \leq 1$ invariant already forbids a second critical section
from being open concurrently regardless of site.

\begin{schema}{PollTickNotifyOk}
  \Delta NotificationState
  \where
  pendingAge' = pendingAge \\
  toolContext = zfalse \\
  siteSuspenders \in inFlight \\
  inFlight' = inFlight \setminus \{ siteSuspenders \} \\
  lastNotifyResult' = notifySent \\
  barStale' = zfalse \\
  session' = session \\
  pendingNotify' = zfalse \\
  notifyLost' = notifyLost \\
  natsConn' = natsConn \\
  kvProvisioned' = kvProvisioned \\
  talkSubscribed' = talkSubscribed \\
  toolContext' = toolContext \\
  readDesc' = readDesc \\
  unreadCount' = unreadCount \\
  talkPending' = talkPending \\
  talkBare' = talkBare \\
  talkQueued' = talkQueued \\
  talkConnected' = talkConnected \\
  wallDesc' = wallDesc \\
  talkDesc' = talkDesc \\
  wallDescChanged' = zfalse \\
  talkDescChanged' = zfalse \\
  queueItems' = queueItems \\
  queueKinds' = queueKinds \\
  queueTexts' = queueTexts \\
  currentIndex' = currentIndex \\
  fileExists' = fileExists \\
  fileUnreadCount' = fileUnreadCount \\
  fileItems' = fileItems \\
  barLine2' = barLine2 \\
  beltPnSnapshot' = beltPnSnapshot \\
  capPnSnapshot' = capPnSnapshot \\
  raceLost' = raceLost \\
  staleWithdrawHit' = staleWithdrawHit
\end{schema}

$PollTickNotifyFail$ is the corrected failure commit: the send fails,
the session reference is cleared exactly as the code above clears it,
the changed flag that triggered the attempt is left set (the client
was not, in fact, told anything), and---the fix---$pendingNotify$ is
raised to record the debt.

\begin{schema}{PollTickNotifyFail}
  \Delta NotificationState
  \where
  pendingAge' = pendingAge \\
  toolContext = zfalse \\
  siteSuspenders \in inFlight \\
  inFlight' = inFlight \setminus \{ siteSuspenders \} \\
  lastNotifyResult' = notifyDropped \\
  barStale' = barStale \\
  session' = sessNone \\
  pendingNotify' = ztrue \\
  notifyLost' = notifyLost \\
  natsConn' = natsConn \\
  kvProvisioned' = kvProvisioned \\
  talkSubscribed' = talkSubscribed \\
  toolContext' = toolContext \\
  readDesc' = readDesc \\
  unreadCount' = unreadCount \\
  talkPending' = talkPending \\
  talkBare' = talkBare \\
  talkQueued' = talkQueued \\
  talkConnected' = talkConnected \\
  wallDesc' = wallDesc \\
  talkDesc' = talkDesc \\
  wallDescChanged' = wallDescChanged \\
  talkDescChanged' = talkDescChanged \\
  queueItems' = queueItems \\
  queueKinds' = queueKinds \\
  queueTexts' = queueTexts \\
  currentIndex' = currentIndex \\
  fileExists' = fileExists \\
  fileUnreadCount' = fileUnreadCount \\
  fileItems' = fileItems \\
  barLine2' = barLine2 \\
  beltPnSnapshot' = beltPnSnapshot \\
  capPnSnapshot' = capPnSnapshot \\
  raceLost' = raceLost \\
  staleWithdrawHit' = staleWithdrawHit
\end{schema}

$PollTickNotifyFailUnrecorded$ is root cause~\#1 itself, kept as the
negative result it now is: identical precondition and identical
$session' = sessNone$, but $pendingNotify$ is left exactly as the
\texttt{return} in the code above leaves it---untouched.
$notifyLost' = ztrue$ records that this has happened, the same device
$staleWithdrawHit$ uses for the talk channel's reordering defect.

\begin{schema}{PollTickNotifyFailUnrecorded}
  \Delta NotificationState
  \where
  pendingAge' = pendingAge \\
  toolContext = zfalse \\
  siteSuspenders \in inFlight \\
  inFlight' = inFlight \setminus \{ siteSuspenders \} \\
  lastNotifyResult' = notifyDropped \\
  barStale' = barStale \\
  session' = sessNone \\
  pendingNotify' = pendingNotify \\
  notifyLost' = ztrue \\
  natsConn' = natsConn \\
  kvProvisioned' = kvProvisioned \\
  talkSubscribed' = talkSubscribed \\
  toolContext' = toolContext \\
  readDesc' = readDesc \\
  unreadCount' = unreadCount \\
  talkPending' = talkPending \\
  talkBare' = talkBare \\
  talkQueued' = talkQueued \\
  talkConnected' = talkConnected \\
  wallDesc' = wallDesc \\
  talkDesc' = talkDesc \\
  wallDescChanged' = wallDescChanged \\
  talkDescChanged' = talkDescChanged \\
  queueItems' = queueItems \\
  queueKinds' = queueKinds \\
  queueTexts' = queueTexts \\
  currentIndex' = currentIndex \\
  fileExists' = fileExists \\
  fileUnreadCount' = fileUnreadCount \\
  fileItems' = fileItems \\
  barLine2' = barLine2 \\
  beltPnSnapshot' = beltPnSnapshot \\
  capPnSnapshot' = capPnSnapshot \\
  raceLost' = raceLost \\
  staleWithdrawHit' = staleWithdrawHit
\end{schema}

\subsection{Recovery: a mid-session reconnect}

$CaptureSession$ (Section~\ref{sec:sessioncapture}) is guarded to the
server's startup window---$natsConn = natsDisconnected$---and cannot
fire again once NATS is up.  But the real
\texttt{\_SessionCaptureMiddleware} runs \texttt{capture\_session()} on
\emph{every} client \texttt{initialize}, not only the first: a client
that reconnects mid-session (a fresh MCP transport, the same NATS
connection underneath) re-arrives here.  $RecaptureSession$ below is
the no-flush case of that second call, exactly mirroring $CaptureSession$'s
own no-flush case: nothing was pending, so the assignment is
synchronous and atomic.

\begin{schema}{RecaptureSession}
  \Delta NotificationState
  \where
  pendingAge' = pendingAge \\
  session = sessNone \\
  natsConn = natsConnected \\
  pendingNotify = zfalse \\
  session' = sessCaptured \\
  natsConn' = natsConn \\
  lastNotifyResult' = lastNotifyResult \\
  wallDescChanged' = wallDescChanged \\
  talkDescChanged' = talkDescChanged \\
  pendingNotify' = pendingNotify \\
  notifyLost' = notifyLost \\
  kvProvisioned' = kvProvisioned \\
  talkSubscribed' = talkSubscribed \\
  toolContext' = toolContext \\
  readDesc' = readDesc \\
  unreadCount' = unreadCount \\
  talkPending' = talkPending \\
  talkBare' = talkBare \\
  talkQueued' = talkQueued \\
  talkConnected' = talkConnected \\
  wallDesc' = wallDesc \\
  talkDesc' = talkDesc \\
  queueItems' = queueItems \\
  queueKinds' = queueKinds \\
  queueTexts' = queueTexts \\
  currentIndex' = currentIndex \\
  fileExists' = fileExists \\
  fileUnreadCount' = fileUnreadCount \\
  fileItems' = fileItems \\
  barLine2' = barLine2 \\
  barStale' = barStale \\
  inFlight' = inFlight \\
  beltPnSnapshot' = beltPnSnapshot \\
  capPnSnapshot' = capPnSnapshot \\
  raceLost' = raceLost \\
  staleWithdrawHit' = staleWithdrawHit
\end{schema}

\subsection{Capture/Reconnect Flush}
\label{sec:capflush}

Both $CaptureSession$ and $RecaptureSession$ call the same
\texttt{capture\_session()}, and when $pendingNotify = ztrue$ that
function does \texttt{await session.send\_tool\_list\_changed()}
before returning---the same shape of critical section as the belt and
suspenders paths, and subject to the same two rounds of findings:
round~1's missing re-arm on failure, and round~2's non-atomicity.
$CaptureSessionFlushBegin$ fires from either entry context (the
disjunction on $natsConn$ is exactly $CaptureSession$'s and
$RecaptureSession$'s two guards); $CaptureSessionFlushOk$ and
$CaptureSessionFlushFail$ are its two commits, sharing one critical
section and one $capPnSnapshot$ because they are, in the running
system, the same function.

\begin{schema}{CaptureSessionFlushBegin}
  \Delta NotificationState
  \where
  pendingAge' = pendingAge \\
  session = sessNone \\
  (natsConn = natsDisconnected \lor natsConn = natsConnected) \\
  pendingNotify = ztrue \\
  siteCapture \notin inFlight \\
  inFlight' = inFlight \cup \{ siteCapture \} \\
  capPnSnapshot' = pendingNotify \\
  session' = session \\
  natsConn' = natsConn \\
  lastNotifyResult' = lastNotifyResult \\
  wallDescChanged' = wallDescChanged \\
  talkDescChanged' = talkDescChanged \\
  pendingNotify' = pendingNotify \\
  notifyLost' = notifyLost \\
  kvProvisioned' = kvProvisioned \\
  talkSubscribed' = talkSubscribed \\
  toolContext' = toolContext \\
  readDesc' = readDesc \\
  unreadCount' = unreadCount \\
  talkPending' = talkPending \\
  talkBare' = talkBare \\
  talkQueued' = talkQueued \\
  talkConnected' = talkConnected \\
  wallDesc' = wallDesc \\
  talkDesc' = talkDesc \\
  queueItems' = queueItems \\
  queueKinds' = queueKinds \\
  queueTexts' = queueTexts \\
  currentIndex' = currentIndex \\
  fileExists' = fileExists \\
  fileUnreadCount' = fileUnreadCount \\
  fileItems' = fileItems \\
  barLine2' = barLine2 \\
  barStale' = barStale \\
  beltPnSnapshot' = beltPnSnapshot \\
  raceLost' = raceLost \\
  staleWithdrawHit' = staleWithdrawHit
\end{schema}

$CaptureSessionFlushOk$: the flush send succeeds, delivering whatever
the tool descriptions currently hold.  $raceLost$ triggers on the same
stale-snapshot test as $NotifyBelt$'s.

\begin{schema}{CaptureSessionFlushOk}
  \Delta NotificationState
  \where
  pendingAge' = pendingAge \\
  session = sessNone \\
  siteCapture \in inFlight \\
  inFlight' = inFlight \setminus \{ siteCapture \} \\
  session' = sessCaptured \\
  natsConn' = natsConn \\
  lastNotifyResult' = notifySent \\
  wallDescChanged' = zfalse \\
  talkDescChanged' = zfalse \\
  pendingNotify' = zfalse \\
  ((pendingNotify \neq capPnSnapshot) \implies raceLost' = ztrue) \\
  ((pendingNotify = capPnSnapshot) \implies raceLost' = raceLost) \\
  notifyLost' = notifyLost \\
  kvProvisioned' = kvProvisioned \\
  talkSubscribed' = talkSubscribed \\
  toolContext' = toolContext \\
  readDesc' = readDesc \\
  unreadCount' = unreadCount \\
  talkPending' = talkPending \\
  talkBare' = talkBare \\
  talkQueued' = talkQueued \\
  talkConnected' = talkConnected \\
  wallDesc' = wallDesc \\
  talkDesc' = talkDesc \\
  queueItems' = queueItems \\
  queueKinds' = queueKinds \\
  queueTexts' = queueTexts \\
  currentIndex' = currentIndex \\
  fileExists' = fileExists \\
  fileUnreadCount' = fileUnreadCount \\
  fileItems' = fileItems \\
  barLine2' = barLine2 \\
  barStale' = barStale \\
  beltPnSnapshot' = beltPnSnapshot \\
  staleWithdrawHit' = staleWithdrawHit
\end{schema}

$CaptureSessionFlushFail$ is round~2's fix at this site: the flush send
fails, the session reference reverts to $sessNone$ (the real code's
exception handler clears \texttt{\_session} here too), and
$pendingNotify$ is re-armed rather than left cleared.
$CaptureSessionFlushFailUnrecorded$ is the dead twin: the real code's
\texttt{\_pending\_notify = False} runs \emph{before} the
\texttt{await}, so an exception handler that does not explicitly
re-arm leaves it cleared---exactly the shape of round~1's suspenders
defect, at a different site.

\begin{schema}{CaptureSessionFlushFail}
  \Delta NotificationState
  \where
  pendingAge' = pendingAge \\
  session = sessNone \\
  siteCapture \in inFlight \\
  inFlight' = inFlight \setminus \{ siteCapture \} \\
  session' = sessNone \\
  natsConn' = natsConn \\
  lastNotifyResult' = notifyDropped \\
  wallDescChanged' = wallDescChanged \\
  talkDescChanged' = talkDescChanged \\
  pendingNotify' = ztrue \\
  notifyLost' = notifyLost \\
  kvProvisioned' = kvProvisioned \\
  talkSubscribed' = talkSubscribed \\
  toolContext' = toolContext \\
  readDesc' = readDesc \\
  unreadCount' = unreadCount \\
  talkPending' = talkPending \\
  talkBare' = talkBare \\
  talkQueued' = talkQueued \\
  talkConnected' = talkConnected \\
  wallDesc' = wallDesc \\
  talkDesc' = talkDesc \\
  queueItems' = queueItems \\
  queueKinds' = queueKinds \\
  queueTexts' = queueTexts \\
  currentIndex' = currentIndex \\
  fileExists' = fileExists \\
  fileUnreadCount' = fileUnreadCount \\
  fileItems' = fileItems \\
  barLine2' = barLine2 \\
  barStale' = barStale \\
  beltPnSnapshot' = beltPnSnapshot \\
  raceLost' = raceLost \\
  staleWithdrawHit' = staleWithdrawHit
\end{schema}

\begin{schema}{CaptureSessionFlushFailUnrecorded}
  \Delta NotificationState
  \where
  pendingAge' = pendingAge \\
  session = sessNone \\
  siteCapture \in inFlight \\
  inFlight' = inFlight \setminus \{ siteCapture \} \\
  session' = sessNone \\
  natsConn' = natsConn \\
  lastNotifyResult' = notifyDropped \\
  wallDescChanged' = wallDescChanged \\
  talkDescChanged' = talkDescChanged \\
  pendingNotify' = zfalse \\
  notifyLost' = ztrue \\
  kvProvisioned' = kvProvisioned \\
  talkSubscribed' = talkSubscribed \\
  toolContext' = toolContext \\
  readDesc' = readDesc \\
  unreadCount' = unreadCount \\
  talkPending' = talkPending \\
  talkBare' = talkBare \\
  talkQueued' = talkQueued \\
  talkConnected' = talkConnected \\
  wallDesc' = wallDesc \\
  talkDesc' = talkDesc \\
  queueItems' = queueItems \\
  queueKinds' = queueKinds \\
  queueTexts' = queueTexts \\
  currentIndex' = currentIndex \\
  fileExists' = fileExists \\
  fileUnreadCount' = fileUnreadCount \\
  fileItems' = fileItems \\
  barLine2' = barLine2 \\
  barStale' = barStale \\
  beltPnSnapshot' = beltPnSnapshot \\
  raceLost' = raceLost \\
  staleWithdrawHit' = staleWithdrawHit
\end{schema}

Root cause~\#1's second consequence, restated precisely now that the
flush is its own critical section: a reconnect that arrives after
$PollTickNotifyFailUnrecorded$ (round~1's defect, $pendingNotify$ still
$zfalse$) finds nothing to flush and takes the no-flush
$RecaptureSession$ path above---restoring $session$ to $sessCaptured$
but delivering nothing, because nothing recorded that anything was
owed.  Once round~1's fix is in place, $pendingNotify = ztrue$ survives
to the reconnect, $CaptureSessionFlushBegin$ fires instead, and the
missed notification is delivered as part of the flush---\emph{unless}
round~2's race clobbers it first, which Section~\ref{sec:race} takes up
next.

\subsection{The counterexample}

Deleting $notifyLost = zfalse$ from $NotificationState$
(Section~\ref{sec:state}) and setting the ProB goal to
$notifyLost = ztrue \land wallDescChanged = ztrue \land talkDescChanged
= zfalse$---isolating the wall channel so the trace is not cluttered by
unrelated talk transitions---finds the four-step trace:

\begin{verbatim}
probcli docs/notification.tex -model_check -noltl -coverage \
  -goal "notifyLost = ztrue & wallDescChanged = ztrue
         & talkDescChanged = zfalse" \
  -p DEFAULT_SETSIZE 2 \
  -card USER 1 -card SOURCEKEY 1 -card SESSIONKEY 1 \
  -card QSLOT 1 -card USLOT 1 -card ASLOT 1 -card PSLOT 1 \
  -p MAX_INITIALISATIONS 16 -p MAX_OPERATIONS 8000
\end{verbatim}

\begin{verbatim}
Found GOAL in state id 1591
Model checking time: 291 ms (304 ms walltime)
States analysed: 423
Transitions fired: 4630
*** COUNTER EXAMPLE FOUND ***
*** TRACE (length=6):
 1: SETUP_CONSTANTS(...)
 2: INITIALISATION(..., natsConn=natsDisconnected, notifyLost=zfalse,
    pendingNotify=zfalse, session=sessNone, ...)
 3: CaptureSession(session=sessCaptured)
 4: EnsureConnected(kvProvisioned=ztrue,natsConn=natsConnected)
 5: KVWallReceive(SOURCEKEY1,MSG2)
 6: PollTickNotifyFailUnrecorded
\end{verbatim}

Step~5 mutates the wall description and sets $wallDescChanged' =
ztrue$; nothing has been sent yet, that is the poller's job on its
next tick.  Step~6 is the suspenders send failing on that very next
tick: $session' = sessNone$, $wallDescChanged$ is left $ztrue$ (the
client was not told), $pendingNotify' = pendingNotify = zfalse$ (the
drop is not recorded), $notifyLost' = ztrue$.  The resulting state is
exactly the one the task asks for: $wallDescChanged = ztrue$ (activity
is server-side correct and waiting), $session = sessNone$ (the
suspenders path is dead), $pendingNotify = zfalse$ (nothing will tell
the next capture to flush), and $natsConn = natsConnected$ throughout
(the connection itself was never at fault).

$RecaptureSession$ is enabled from this state---its guard $session =
sessNone \land natsConn = natsConnected \land pendingNotify = zfalse$
holds (round~1's defect never sets $pendingNotify$)---and confirms the
second half of the finding without needing a further ProB search: by
inspection of $RecaptureSession$'s guard and body
(Section~\ref{sec:capflush}), firing it sets $session' = sessCaptured$
while leaving $wallDescChanged'$, $talkDescChanged'$, and
$lastNotifyResult'$ all unchanged.  The reconnect does not recover the
missed notification, because nothing recorded that one was owed.

\begin{verbatim}
probcli docs/notification.tex -model_check -noltl -coverage \
  -p DEFAULT_SETSIZE 2 \
  -card USER 1 -card SOURCEKEY 1 -card SESSIONKEY 1 \
  -card QSLOT 1 -card USLOT 1 -card ASLOT 1 -card PSLOT 1 \
  -p MAX_INITIALISATIONS 16 -p MAX_OPERATIONS 8000
\end{verbatim}

\paragraph{What this does and does not establish.}  Round~1's own
exhaustive run at this scope---$1{,}086{,}889$ states,
$invariant\_violated{:}0$, $deadlocked{:}0$, $open{:}0$,
$PollTickNotifyFailUnrecorded$ and $NotifyBeltStale$ (the earlier,
single-site name for the belt defect, since retired---see
Section~\ref{sec:beltpath}) reported uncovered---was computed \emph{before}
round~2 added $CaptureSessionFlushBegin$, $NotifyBeltBegin$,
$PollTickNotifyBegin$, and the $inFlight$/snapshot machinery.  Re-running
the same exhaustive, no-goal check against the extended schema at the
same floor scope did not terminate in a reasonable time or memory
budget---the reachable space, thirty-odd talk/wall/read marker
components combined with the new critical-section splitting, grew past
$6.5$ million states and was still climbing after twenty minutes when
the run was stopped.  None of that growth comes from the property being
checked: $notifyLost = zfalse$ depends on five variables
($pendingNotify$, $inFlight$, the two snapshots, itself), not on the
talk/wall/read machinery it happened to be checked alongside.
\texttt{notification-race.tex}, cited in full in
Section~\ref{sec:race}, restates the property (and $raceLost = zfalse$
alongside it) against exactly those five variables in their own
eleven-operation machine, and both hold exhaustively there in
$23$ states and $33$~ms, with $PollTickNotifyFailUnrecorded$,
$NotifyBeltSendFailUnrecorded$, and $CaptureSessionFlushFailUnrecorded$
all reported uncovered---the same proof shape as $WithdrawStale$
(Section~\ref{sec:markerfix2}), and the authority for the claim that
follows.

This is a \emph{safety} property: it says the failure-accounting
obligation is never skipped, in any reachable state.  It is
deliberately not phrased as the CTL liveness formula
$\mathrm{AG}\,((wallDescChanged = ztrue \lor talkDescChanged = ztrue)
\implies \mathrm{EF}\,(wallDescChanged = zfalse \land talkDescChanged =
zfalse))$ that closes the talk channel's own defect in
Section~\ref{sec:markerfix2}, because that formula would report
\textsc{true} in \emph{both} the buggy and the fixed model here: even
though $NotifyBelt$ can now fail too, its success commit remains
enabled whenever $toolContext = ztrue$ and a belt call is in flight, so
ProB's existential path quantifier $\mathrm{EF}$ always finds a
recovering trace through some future, unrelated tool call succeeding,
whether or not the suspenders bookkeeping was correct.  That is not a
false proof, but it is not the property this bug is about, and stating
it as though it were would overclaim what the check
shows---the {\fuzz} discipline is to report exactly what a green check
establishes, not what would be reassuring for it to establish.  The
talk channel's liveness proof works because $ExpirePendingInvite$ is
driven by a clock ($AgeTick$, $pendingAge$) guaranteed to advance on
its own; nothing plays that role for the suspenders path once $session$
is lost.  The operations that restore it---$CaptureSession$,
$RecaptureSession$, $CaptureSessionFlushOk$, $NotifyBelt$---are all
environment-triggered (server startup, a client reconnect, a fresh tool
call) rather than time-driven, so no purely background, poller-only
sequence of transitions ever clears a stranded changed flag.  That
structural fact is visible by inspection of the schema set above: no
operation whose guard depends only on $toolContext = zfalse$ sets
$session' = sessCaptured$ from $session = sessNone$.  The safety
property $notifyLost = zfalse$ is the one this specification actually
proves, and it is the one the fix needs: once every suspenders failure
is recorded, the next capture point---whichever arrives
first---discharges it, by inspection of $CaptureSession$'s,
$RecaptureSession$'s, and $CaptureSessionFlushOk$'s guards.

%%----------------------------------------------------------------------
\section{The Race: Three Critical Sections, One Shared Flag}
\label{sec:race}
%%----------------------------------------------------------------------

Every operation above that writes $pendingNotify$ was, until this
section, modelled as a single atomic step: read the old value, decide
the new one, send, write---all in one transition.  The real code is not
atomic there.  Each of the three sites does
\texttt{read-or-decide; await send(...); write}, and Python's
cooperative scheduler can run any other coroutine during the
\texttt{await}.  $CriticalSite$ and $inFlight$ (Section~\ref{sec:state})
exist to make that gap a reachable state: $Begin$ at each site records
that the read/decide has happened and the write is still outstanding;
$Commit$ performs the write.  Between a site's $Begin$ and its $Commit$,
any \emph{other} site's $Begin$ can fire---unless something forbids it.

\paragraph{The race, reproduced.}  Take $\# inFlight \leq 1$ out of
$NotificationState$ (Section~\ref{sec:state}) and set the ProB goal to
$raceLost = ztrue$:

\begin{verbatim}
probcli docs/notification.tex -model_check -noltl -coverage \
  -goal "raceLost = ztrue & natsConn = natsConnected
         & pendingNotify = zfalse" \
  -p DEFAULT_SETSIZE 2 \
  -card USER 1 -card SOURCEKEY 1 -card SESSIONKEY 1 \
  -card QSLOT 1 -card USLOT 1 -card ASLOT 1 -card PSLOT 1 \
  -p MAX_INITIALISATIONS 16 -p MAX_OPERATIONS 8000
\end{verbatim}

\begin{verbatim}
Found GOAL in state id 11897
States analysed: 5985   Transitions fired: 87919
*** COUNTER EXAMPLE FOUND ***
*** TRACE (length=12):
 1: SETUP_CONSTANTS(...)
 2: INITIALISATION(..., session=sessNone, pendingNotify=zfalse,
    inFlight={}, raceLost=zfalse, ...)
 3: EnterToolContext
 4: NotifyBeltBegin
 5: NotifyBeltSendFail
 6: SyncUnreadFile(0)
 7: StatusBarRefresh(barLine2=MSG2)
 8: CaptureSessionFlushBegin
 9: NotifyBeltBegin
 10: CaptureSessionFlushOk(inFlight={siteBelt},lastNotifyResult=notifySent,
     pendingNotify=zfalse,session=sessCaptured)
 11: EnsureConnected(kvProvisioned=ztrue,natsConn=natsConnected)
 12: NotifyBelt(barStale=zfalse,inFlight={},raceLost=ztrue)
\end{verbatim}

ProB's shortest witness is a belt/capture collision, not the
belt/suspenders one sketched informally above---but it is the same
mechanism, and it is worth walking precisely rather than
re-describing the informal sketch as though it were what was found.
Step~4--5: a first belt call's send fails; its commit re-arms
$pendingNotify' = ztrue$.  Step~8: capture's flush begins, reading
that re-armed $ztrue$ and snapshotting $capPnSnapshot = ztrue$.
Step~9: \emph{before capture commits}, a \emph{second} belt call
begins---only reachable because the mutex invariant is absent, so
$siteBelt$ can re-enter $inFlight$ once its first occupant released it
at step~5---and this second call snapshots $beltPnSnapshot = ztrue$
too, reading the same still-outstanding debt.  Step~10: capture's
flush succeeds first, correctly delivering the notification and
clearing $pendingNotify' = zfalse$.  Step~12: the second belt call's
commit, made on a decision formed back at step~9 when $pendingNotify$
was still $ztrue$, now finds $pendingNotify = zfalse \neq
beltPnSnapshot = ztrue$---the value it decided against has already
changed---and $raceLost' = ztrue$ fires.

In this particular interleaving the final value of $pendingNotify$
happens to still be correct ($zfalse$, and it was genuinely
delivered---by capture, at step~10): the harm here is redundant,
order-dependent work, not a lost debt.  $raceLost$ is deliberately the
more conservative check---any commit whose snapshot is stale at commit
time, not only ones that provably corrupt the final value---because
the \emph{class} of behaviour is unsafe regardless of whether a given
interleaving happens to land safely: change step~10's outcome to a
failure instead of a success (a second, independent send failure,
fully reachable by the same construction) and the identical
stale-snapshot commit at step~12 clobbers a live re-arm exactly as the
belt/suspenders sketch described.  $raceLost = ztrue$ is the witness
for that whole class, not only for its most damaging member.

\paragraph{The mutex.}  $\# inFlight \leq 1$ is the Z-level statement of
a lock held for the duration of each site's critical section: at most
one of $siteSuspenders$, $siteBelt$, $siteCapture$ may have a write
outstanding at any reachable state, regardless of which two sites
collide.  Restoring it makes every $Begin$ whose site is not already
the sole occupant of $inFlight$ infeasible while another site is
mid-flight---not by adding a new guard to each $Begin$ schema, but
because any reachable state with two sites in $inFlight$ would violate
the promoted invariant, so ProB never visits one.  Under the mutex,
step~9 above cannot happen between steps~8 and~10: the second
$NotifyBeltBegin$'s $siteBelt \notin inFlight$ guard is satisfiable
(belt's own first occupancy already released at step~5), but the
post-state $inFlight' = \{siteCapture, siteBelt\}$ has
$\# inFlight' = 2$, contradicting the invariant, so the operation is
enabled in no reachable state while $siteCapture$ is still held.  The
second belt call must wait for capture's commit (step~10) to release
the lock first; by the time its $Begin$ can fire, its snapshot
$beltPnSnapshot' = pendingNotify$ reads the value \emph{capture just
delivered}, $zfalse$, so its own decision is now ``nothing to
flush''---it commits a genuine no-op rather than a stale, redundant
clear, and $raceLost$ never fires.  The same argument applies verbatim
to a suspenders/belt or suspenders/capture pairing: the mutex forbids
\emph{any} two of the three sites having a write outstanding at once,
so whichever pair a given execution would have collided on, the second
to attempt $Begin$ waits and reads the first's committed result
instead of a value it will overwrite.

Running that check---$notifyLost = zfalse$ and $\# inFlight \leq 1$
both restored---against the full $NotificationState$ at this floor
scope is exactly the run that did not terminate
(Section~\ref{sec:sessionloss}'s ``What this does and does not
establish''): $6.5$ million states and climbing after twenty minutes,
stopped rather than left to find out how much further it would grow.
The reachable-space explosion is the talk/wall/read marker machinery's,
not the race's---nothing about $\# inFlight \leq 1$ or $raceLost$
depends on $talkPending$, $queueItems$, or any of the other two dozen
components $NotificationState$ carries for unrelated channels.
\texttt{notification-race.tex} isolates the five variables the mutex
and the race actually touch as their own machine and checks the
identical property there:

\begin{verbatim}
probcli docs/notification-race.tex -model_check -noltl -coverage \
  -p DEFAULT_SETSIZE 2 -p MAX_OPERATIONS 2000
\end{verbatim}

$23$ states, complete, in $33$~ms: $invariant\_violated{:}0$,
$deadlocked{:}0$, $open{:}0$.  $PollTickNotifyFailUnrecorded$,
$NotifyBeltSendFailUnrecorded$, and $CaptureSessionFlushFailUnrecorded$
are reported uncovered there (round~1's finding, now proved at all
three sites in the same run); no operation is reported uncovered
\emph{because of} the mutex---every $Begin$/$Commit$ pair is exercised
at least once, just never with a second site's $Begin$ landing while
the first is still mid-flight.  $raceLost = zfalse$ and $notifyLost =
zfalse$ both hold across the whole reachable space of that isolated
machine.  The schemas there are verbatim restrictions of the ones
above---$SusBegin$/$SusCommitSet$ to $PollTickNotifyBegin$/$Fail$,
$BeltBegin$/$BeltCommitClear$/$BeltCommitSet$ to
$NotifyBeltBegin$/$NotifyBelt$/$NotifyBeltSendFail$, and
$CapBegin$/$CapCommitClear$/$CapCommitSet$ to
$CaptureSessionFlushBegin$/$Ok$/$Fail$---so the isolated proof
transfers: the omitted variables cannot make the mutex insufficient,
because none of them appear in any guard or effect that touches
$pendingNotify$, $inFlight$, or the two snapshots.

\paragraph{Locking scope, for the implementation.}  The mutex proved
above is a single lock shared by all three sites---not three separate
locks, because the race is precisely between \emph{different} sites,
and three independent locks would each protect their own site from
itself while leaving every cross-site interleaving open.  A single
\texttt{asyncio.Lock()} at module scope in \texttt{\_descriptions.py},
held across exactly the following region at each site, is what the
proof requires:

\begin{itemize}
  \item \textbf{Suspenders} (\texttt{notify\_tool\_list\_changed()},
    suspenders branch): acquire before the
    \texttt{if \_session is not None:} read; hold across
    \texttt{await \_session.send\_tool\_list\_changed()}; release only
    after \emph{both} \texttt{\_session} and \texttt{\_pending\_notify}
    have been written in the \texttt{except} branch (or, on success,
    after the function returns).  The pre-session branch
    (\texttt{\_pending\_notify = True}) is a plain synchronous write
    with no \texttt{await} in it and does not need the lock on its own
    account, but taking it anyway is harmless and keeps one acquisition
    point per call.
  \item \textbf{Belt} (\texttt{notify\_tool\_list\_changed()}, belt
    branch): acquire before \texttt{ctx = get\_context()}; hold across
    \texttt{await ctx.send\_notification(...)}; release only after
    \texttt{\_session = ctx.session} on success, or after the (new,
    round~2) \texttt{\_pending\_notify = True} re-arm on a caught
    non-\texttt{RuntimeError} failure.  Do \emph{not} release between
    the send and the session/pending-notify write---that gap is exactly
    what $beltPnSnapshot$ models, and releasing early reopens the race
    the proof just closed.
  \item \textbf{Capture/reconnect flush} (\texttt{capture\_session()}):
    acquire before the \texttt{if not \_pending\_notify: return} read
    (a session that has nothing to flush still needs the lock to
    observe a consistent $pendingNotify$, matching $CaptureSessionFlushBegin$'s
    read of it); hold across \texttt{await
    session.send\_tool\_list\_changed()}; release only after
    \texttt{\_pending\_notify} and \texttt{\_session} are both settled
    in the \texttt{except} branch (round~2: set
    \texttt{\_pending\_notify = True} there, matching
    $CaptureSessionFlushFail$), or after the function returns on
    success.  \texttt{\_session = session}, the unconditional first line
    of \texttt{capture\_session()}, happens before the
    $pendingNotify$ check and is not itself part of this critical
    section---it has no failure mode and nothing else reads or writes
    $pendingNotify$ around it.
\end{itemize}

The scope is the same shape at all three sites: \emph{acquire before
the last read of \texttt{\_pending\_notify}, release after the last
write of \texttt{\_pending\_notify} (and, where the site also owns it,
\texttt{\_session})---with the \texttt{await} in between} held under
the lock throughout.  A lock acquired only around the write, or
released before the write, does not close the race: $beltPnSnapshot$
and $capPnSnapshot$ are taken at $Begin$ precisely because the proof
depends on the \emph{read} being inside the same held region as the
eventual write, not just the write being atomic on its own.

\subsection{Status Bar Refresh}

Claude~Code polls the status line command on its own schedule.
Reads the filesystem status file and updates the display.
This is the slow path---no MCP involvement.

\begin{schema}{StatusBarRefresh}
  \Delta NotificationState
  \where
  pendingAge' = pendingAge \\
  (fileExists = ztrue \land \\
  \quad~ fileItems \neq \langle \rangle \implies \\
  \quad~ barLine2' \neq idleMarker) \\
  (fileExists = zfalse \implies
  barLine2' = idleMarker) \\
  natsConn' = natsConn \\
  kvProvisioned' = kvProvisioned \\
  talkSubscribed' = talkSubscribed \\
  session' = session \\
  toolContext' = toolContext \\
  readDesc' = readDesc \\
  unreadCount' = unreadCount \\
  talkPending' = talkPending \\
  talkBare' = talkBare \\
  talkQueued' = talkQueued \\
  talkConnected' = talkConnected \\
  wallDesc' = wallDesc \\
  talkDesc' = talkDesc \\
  wallDescChanged' = wallDescChanged \\
  talkDescChanged' = talkDescChanged \\
  queueItems' = queueItems \\
  queueKinds' = queueKinds \\
  queueTexts' = queueTexts \\
  currentIndex' = currentIndex \\
  fileExists' = fileExists \\
  fileUnreadCount' = fileUnreadCount \\
  fileItems' = fileItems \\
  barStale' = barStale \\
  lastNotifyResult' = lastNotifyResult \\
  pendingNotify' = pendingNotify \\
  notifyLost' = notifyLost \\
  inFlight' = inFlight \\
  beltPnSnapshot' = beltPnSnapshot \\
  capPnSnapshot' = capPnSnapshot \\
  raceLost' = raceLost \\
  staleWithdrawHit' = staleWithdrawHit
\end{schema}

%%----------------------------------------------------------------------
\section{Display Queue Operations}
%%----------------------------------------------------------------------

\subsection{Advance If Due}

Rotation: wall items cycle (persistent), talk items are
ephemeral (removed after one display turn).  The $currentIndex + 1$
expression indexes into the 1-based sequence.

Talk items at the current position are removed (ephemeral).
Wall items cycle: the cursor advances, but the item stays in the
queue.  With a single wall item, the cursor stays at zero.

\begin{schema}{AdvanceTalk}
  \Delta NotificationState
  \where
  pendingAge' = pendingAge \\
  queueItems \neq \langle \rangle \\
  queueKinds(queueItems(currentIndex + 1)) = itemTalk \\
  \# queueItems' = \# queueItems - 1 \\
  queueKinds' = \\
  \quad~ \{ queueItems(currentIndex + 1) \} \ndres queueKinds \\
  queueTexts' = \\
  \quad~ \{ queueItems(currentIndex + 1) \} \ndres queueTexts \\
  \ran queueItems' = \ran queueItems \setminus \\
  \quad~ \{ queueItems(currentIndex + 1) \} \\
  natsConn' = natsConn \\
  kvProvisioned' = kvProvisioned \\
  talkSubscribed' = talkSubscribed \\
  session' = session \\
  toolContext' = toolContext \\
  readDesc' = readDesc \\
  unreadCount' = unreadCount \\
  talkPending' = talkPending \\
  talkBare' = talkBare \\
  talkQueued' = talkQueued \\
  talkConnected' = talkConnected \\
  wallDesc' = wallDesc \\
  talkDesc' = talkDesc \\
  wallDescChanged' = wallDescChanged \\
  talkDescChanged' = talkDescChanged \\
  fileExists' = fileExists \\
  fileUnreadCount' = fileUnreadCount \\
  fileItems' = fileItems \\
  barLine2' = barLine2 \\
  barStale' = barStale \\
  lastNotifyResult' = lastNotifyResult \\
  pendingNotify' = pendingNotify \\
  notifyLost' = notifyLost \\
  inFlight' = inFlight \\
  beltPnSnapshot' = beltPnSnapshot \\
  capPnSnapshot' = capPnSnapshot \\
  raceLost' = raceLost \\
  staleWithdrawHit' = staleWithdrawHit
\end{schema}

\begin{schema}{AdvanceWall}
  \Delta NotificationState
  \where
  pendingAge' = pendingAge \\
  queueItems \neq \langle \rangle \\
  queueKinds(queueItems(currentIndex + 1)) = itemWall \\
  queueItems' = queueItems \\
  queueKinds' = queueKinds \\
  queueTexts' = queueTexts \\
  (\# queueItems > 1 \implies \\
  \quad~ currentIndex' = (currentIndex + 1) \mod \\
  \quad~ \quad~ \# queueItems) \\
  (\# queueItems = 1 \implies
  currentIndex' = 0) \\
  natsConn' = natsConn \\
  kvProvisioned' = kvProvisioned \\
  talkSubscribed' = talkSubscribed \\
  session' = session \\
  toolContext' = toolContext \\
  readDesc' = readDesc \\
  unreadCount' = unreadCount \\
  talkPending' = talkPending \\
  talkBare' = talkBare \\
  talkQueued' = talkQueued \\
  talkConnected' = talkConnected \\
  wallDesc' = wallDesc \\
  talkDesc' = talkDesc \\
  wallDescChanged' = wallDescChanged \\
  talkDescChanged' = talkDescChanged \\
  fileExists' = fileExists \\
  fileUnreadCount' = fileUnreadCount \\
  fileItems' = fileItems \\
  barLine2' = barLine2 \\
  barStale' = barStale \\
  lastNotifyResult' = lastNotifyResult \\
  pendingNotify' = pendingNotify \\
  notifyLost' = notifyLost \\
  inFlight' = inFlight \\
  beltPnSnapshot' = beltPnSnapshot \\
  capPnSnapshot' = capPnSnapshot \\
  raceLost' = raceLost \\
  staleWithdrawHit' = staleWithdrawHit
\end{schema}

%%----------------------------------------------------------------------
\section{Tool-Description Marker Channels}
\label{sec:marker}
%%----------------------------------------------------------------------

The MCP server surfaces incoming activity by mutating a tool's
\emph{description} and firing \texttt{tools/list\_changed}: the agent
notices the changed text on its next turn.  Two such channels are
modelled here.  The \texttt{read\_messages} channel is the reference: it
is known to behave correctly.  The \texttt{talk} channel is the one under
examination.

The governing property, developed in Section~\ref{sec:markerprops}, is
that for each channel the marker is present exactly when the underlying
activity exists:
\[
  desc \neq default \iff activity.
\]
This section gives the operations of both channels.  It deliberately
gives them first \emph{without} the reconciling transitions that make the
property hold, so that Section~\ref{sec:markerprops} can exhibit the
defect as a reachable state before Section~\ref{sec:markerfix} repairs
it.

\subsection{Read Channel (the self-clearing reference)}

Mail arrives and raises the unread count; the description gains its
marker.  Reading mail returns the count to zero and the description
reverts.  Because the description is driven by a count that
\emph{decreases}, this channel clears itself: there is no way to leave
the marker set once the count reaches zero.

\begin{schema}{MailArrive}
  \Delta NotificationState
  \where
  pendingAge' = pendingAge \\
  unreadCount < maxUnreadCount \\
  unreadCount' = unreadCount + 1 \\
  readDesc' \neq readDescDefault \\
  natsConn' = natsConn \\
  kvProvisioned' = kvProvisioned \\
  talkSubscribed' = talkSubscribed \\
  session' = session \\
  toolContext' = toolContext \\
  wallDesc' = wallDesc \\
  talkDesc' = talkDesc \\
  wallDescChanged' = wallDescChanged \\
  talkDescChanged' = talkDescChanged \\
  talkPending' = talkPending \\
  talkBare' = talkBare \\
  talkQueued' = talkQueued \\
  talkConnected' = talkConnected \\
  queueItems' = queueItems \\
  queueKinds' = queueKinds \\
  queueTexts' = queueTexts \\
  currentIndex' = currentIndex \\
  fileExists' = fileExists \\
  fileUnreadCount' = fileUnreadCount \\
  fileItems' = fileItems \\
  barLine2' = barLine2 \\
  barStale' = barStale \\
  lastNotifyResult' = lastNotifyResult \\
  pendingNotify' = pendingNotify \\
  notifyLost' = notifyLost \\
  inFlight' = inFlight \\
  beltPnSnapshot' = beltPnSnapshot \\
  capPnSnapshot' = capPnSnapshot \\
  raceLost' = raceLost \\
  staleWithdrawHit' = staleWithdrawHit
\end{schema}

\begin{schema}{ReadMail}
  \Delta NotificationState
  \where
  pendingAge' = pendingAge \\
  unreadCount > 0 \\
  unreadCount' = 0 \\
  readDesc' = readDescDefault \\
  natsConn' = natsConn \\
  kvProvisioned' = kvProvisioned \\
  talkSubscribed' = talkSubscribed \\
  session' = session \\
  toolContext' = toolContext \\
  wallDesc' = wallDesc \\
  talkDesc' = talkDesc \\
  wallDescChanged' = wallDescChanged \\
  talkDescChanged' = talkDescChanged \\
  talkPending' = talkPending \\
  talkBare' = talkBare \\
  talkQueued' = talkQueued \\
  talkConnected' = talkConnected \\
  queueItems' = queueItems \\
  queueKinds' = queueKinds \\
  queueTexts' = queueTexts \\
  currentIndex' = currentIndex \\
  fileExists' = fileExists \\
  fileUnreadCount' = fileUnreadCount \\
  fileItems' = fileItems \\
  barLine2' = barLine2 \\
  barStale' = barStale \\
  lastNotifyResult' = lastNotifyResult \\
  pendingNotify' = pendingNotify \\
  notifyLost' = notifyLost \\
  inFlight' = inFlight \\
  beltPnSnapshot' = beltPnSnapshot \\
  capPnSnapshot' = capPnSnapshot \\
  raceLost' = raceLost \\
  staleWithdrawHit' = staleWithdrawHit
\end{schema}

\subsection{Talk Channel (invite, message, drain)}

An invite arrives and is recorded in $talkPending$, retaining the
inviter's session key; the description gains its \texttt{[TALK]} marker.
A talk message arrives and raises $talkQueued$.  Both set the marker.

$TalkInviteArrive$ fires when the invite either supersedes one already
recorded for that user ($inviter? \in \dom talkPending$, no growth) or
there is room for a new entry ($\# talkPending + \# talkBare <
maxPending$).  This is the exact analogue of $ReceiveNotification$'s
$queueLen < maxQueueLen$ guard in $talk.tex$: it keeps the pending set
within its bound.  When a \emph{new} inviter arrives and the set is
already full, $TalkInviteArriveOverflow$ takes over, evicting one
existing entry before inserting---the drop-to-make-room counterpart of
$ReceiveOverflow$.

\begin{schema}{TalkInviteArrive}
  \Delta NotificationState \\
  inviter? : USER \\
  inviterKey? : SESSIONKEY
  \where
  pendingAge' = pendingAge \\
  natsConn = natsConnected \\
  inviter? \notin talkBare \\
  (inviter? \in \dom talkPending \lor
  \# talkPending + \# talkBare < maxPending) \\
  talkPending' = talkPending \oplus \{ inviter? \mapsto inviterKey? \} \\
  talkDesc' \neq talkDescDefault \\
  talkDescChanged' = ztrue \\
  natsConn' = natsConn \\
  kvProvisioned' = kvProvisioned \\
  talkSubscribed' = talkSubscribed \\
  session' = session \\
  toolContext' = toolContext \\
  wallDesc' = wallDesc \\
  readDesc' = readDesc \\
  wallDescChanged' = wallDescChanged \\
  unreadCount' = unreadCount \\
  talkBare' = talkBare \\
  talkQueued' = talkQueued \\
  talkConnected' = talkConnected \\
  queueItems' = queueItems \\
  queueKinds' = queueKinds \\
  queueTexts' = queueTexts \\
  currentIndex' = currentIndex \\
  fileExists' = fileExists \\
  fileUnreadCount' = fileUnreadCount \\
  fileItems' = fileItems \\
  barLine2' = barLine2 \\
  barStale' = barStale \\
  lastNotifyResult' = lastNotifyResult \\
  pendingNotify' = pendingNotify \\
  notifyLost' = notifyLost \\
  inFlight' = inFlight \\
  beltPnSnapshot' = beltPnSnapshot \\
  capPnSnapshot' = capPnSnapshot \\
  raceLost' = raceLost \\
  staleWithdrawHit' = staleWithdrawHit
\end{schema}

The pending set is bounded by $maxPending$ ($\# PSLOT$), exactly as the
message queue is bounded by $maxQueueSize$.  When a \emph{new} inviter
arrives and the set is full, $TalkInviteArriveOverflow$ evicts one
existing entry---the oldest by arrival in the implementation---and
inserts the new one, so the bound is held and no stream of forged
inviters can exhaust memory.  The model records the eviction only as
the removal of \emph{some} pending entry $evicted?$: $talkPending$ is an
unordered function, so arrival order is not expressible here, and
``oldest by arrival'' is the obligation the implementation discharges.
What the model proves is the invariant---$\# talkPending + \# talkBare
\leq maxPending$ across the whole reachable space---not the choice of
victim.  This is the pending-set mirror of $ReceiveOverflow$ in
$talk.tex$, where a full queue drops its head before appending.

\begin{schema}{TalkInviteArriveOverflow}
  \Delta NotificationState \\
  inviter? : USER \\
  inviterKey? : SESSIONKEY \\
  evicted? : USER
  \where
  pendingAge' = pendingAge \\
  natsConn = natsConnected \\
  inviter? \notin talkBare \\
  inviter? \notin \dom talkPending \\
  \# talkPending + \# talkBare = maxPending \\
  evicted? \in \dom talkPending \\
  talkPending' =
  (\{ evicted? \} \ndres talkPending) \oplus
  \{ inviter? \mapsto inviterKey? \} \\
  talkDesc' \neq talkDescDefault \\
  talkDescChanged' = ztrue \\
  natsConn' = natsConn \\
  kvProvisioned' = kvProvisioned \\
  talkSubscribed' = talkSubscribed \\
  session' = session \\
  toolContext' = toolContext \\
  wallDesc' = wallDesc \\
  readDesc' = readDesc \\
  wallDescChanged' = wallDescChanged \\
  unreadCount' = unreadCount \\
  talkBare' = talkBare \\
  talkQueued' = talkQueued \\
  talkConnected' = talkConnected \\
  queueItems' = queueItems \\
  queueKinds' = queueKinds \\
  queueTexts' = queueTexts \\
  currentIndex' = currentIndex \\
  fileExists' = fileExists \\
  fileUnreadCount' = fileUnreadCount \\
  fileItems' = fileItems \\
  barLine2' = barLine2 \\
  barStale' = barStale \\
  lastNotifyResult' = lastNotifyResult \\
  pendingNotify' = pendingNotify \\
  notifyLost' = notifyLost \\
  inFlight' = inFlight \\
  beltPnSnapshot' = beltPnSnapshot \\
  capPnSnapshot' = capPnSnapshot \\
  raceLost' = raceLost \\
  staleWithdrawHit' = staleWithdrawHit
\end{schema}

\begin{schema}{TalkMessageArrive}
  \Delta NotificationState
  \where
  pendingAge' = pendingAge \\
  natsConn = natsConnected \\
  talkQueued < maxQueueSize \\
  talkQueued' = talkQueued + 1 \\
  talkDesc' \neq talkDescDefault \\
  talkDescChanged' = ztrue \\
  natsConn' = natsConn \\
  kvProvisioned' = kvProvisioned \\
  talkSubscribed' = talkSubscribed \\
  session' = session \\
  toolContext' = toolContext \\
  wallDesc' = wallDesc \\
  readDesc' = readDesc \\
  wallDescChanged' = wallDescChanged \\
  unreadCount' = unreadCount \\
  talkPending' = talkPending \\
  talkBare' = talkBare \\
  talkConnected' = talkConnected \\
  queueItems' = queueItems \\
  queueKinds' = queueKinds \\
  queueTexts' = queueTexts \\
  currentIndex' = currentIndex \\
  fileExists' = fileExists \\
  fileUnreadCount' = fileUnreadCount \\
  fileItems' = fileItems \\
  barLine2' = barLine2 \\
  barStale' = barStale \\
  lastNotifyResult' = lastNotifyResult \\
  pendingNotify' = pendingNotify \\
  notifyLost' = notifyLost \\
  inFlight' = inFlight \\
  beltPnSnapshot' = beltPnSnapshot \\
  capPnSnapshot' = capPnSnapshot \\
  raceLost' = raceLost \\
  staleWithdrawHit' = staleWithdrawHit
\end{schema}

The defect (biff-9la, observed live).  The drain that ends a talk---the
aggregate of accepting, exchanging, and hanging up---empties every source
of talk activity, but leaves $talkDesc$ untouched.  Nothing reverts the
marker; there is no withdrawal frame and no expiry.  $TalkDrainStale$
models exactly this: it clears $talkPending$, $talkBare$, $talkQueued$
and $talkConnected$, yet preserves $talkDesc$.

\begin{schema}{TalkDrainStale}
  \Delta NotificationState
  \where
  pendingAge' = 0 \\
  (\dom talkPending \neq \emptyset \lor talkBare \neq \emptyset \lor
  talkQueued > 0 \lor talkConnected = ztrue) \\
  talkPending' = \emptyset \\
  talkBare' = \emptyset \\
  talkQueued' = 0 \\
  talkConnected' = zfalse \\
  talkDesc' = talkDesc \\
  talkDescChanged' = talkDescChanged \\
  natsConn' = natsConn \\
  kvProvisioned' = kvProvisioned \\
  talkSubscribed' = talkSubscribed \\
  session' = session \\
  toolContext' = toolContext \\
  wallDesc' = wallDesc \\
  readDesc' = readDesc \\
  wallDescChanged' = wallDescChanged \\
  unreadCount' = unreadCount \\
  queueItems' = queueItems \\
  queueKinds' = queueKinds \\
  queueTexts' = queueTexts \\
  currentIndex' = currentIndex \\
  fileExists' = fileExists \\
  fileUnreadCount' = fileUnreadCount \\
  fileItems' = fileItems \\
  barLine2' = barLine2 \\
  barStale' = barStale \\
  lastNotifyResult' = lastNotifyResult \\
  pendingNotify' = pendingNotify \\
  notifyLost' = notifyLost \\
  inFlight' = inFlight \\
  beltPnSnapshot' = beltPnSnapshot \\
  capPnSnapshot' = capPnSnapshot \\
  raceLost' = raceLost \\
  staleWithdrawHit' = staleWithdrawHit
\end{schema}

The coupled defect.  A second, independent fault records an invite
\emph{without} its session key---the tty is dropped on the way into the
pending representation.  $TalkInviteArriveBare$ puts the inviter into
$talkBare$ and raises the marker; the resulting accept hint reads a bare
\texttt{talk @user}, which fails at the prompt (``Talk needs a specific
session'').

\begin{schema}{TalkInviteArriveBare}
  \Delta NotificationState \\
  inviter? : USER
  \where
  pendingAge' = pendingAge \\
  natsConn = natsConnected \\
  inviter? \notin \dom talkPending \\
  talkBare' = talkBare \cup \{ inviter? \} \\
  talkDesc' \neq talkDescDefault \\
  talkDescChanged' = ztrue \\
  natsConn' = natsConn \\
  kvProvisioned' = kvProvisioned \\
  talkSubscribed' = talkSubscribed \\
  session' = session \\
  toolContext' = toolContext \\
  wallDesc' = wallDesc \\
  readDesc' = readDesc \\
  wallDescChanged' = wallDescChanged \\
  unreadCount' = unreadCount \\
  talkPending' = talkPending \\
  talkQueued' = talkQueued \\
  talkConnected' = talkConnected \\
  queueItems' = queueItems \\
  queueKinds' = queueKinds \\
  queueTexts' = queueTexts \\
  currentIndex' = currentIndex \\
  fileExists' = fileExists \\
  fileUnreadCount' = fileUnreadCount \\
  fileItems' = fileItems \\
  barLine2' = barLine2 \\
  barStale' = barStale \\
  lastNotifyResult' = lastNotifyResult \\
  pendingNotify' = pendingNotify \\
  notifyLost' = notifyLost \\
  inFlight' = inFlight \\
  beltPnSnapshot' = beltPnSnapshot \\
  capPnSnapshot' = capPnSnapshot \\
  raceLost' = raceLost \\
  staleWithdrawHit' = staleWithdrawHit
\end{schema}

%%----------------------------------------------------------------------
\section{Marker Consistency: the property and the defect}
\label{sec:markerprops}
%%----------------------------------------------------------------------

Write $talkActivity$ for the disjunction of the three talk activity
sources and $readActivity$ for a positive unread count:
\begin{eqnarray*}
  talkActivity & \defs & \dom talkPending \neq \emptyset \lor
    talkBare \neq \emptyset \lor talkQueued > 0 \lor
    talkConnected = ztrue \\
  readActivity & \defs & unreadCount > 0.
\end{eqnarray*}

The intended invariant is that each marker is present exactly when its
activity exists, and that a talk marker always carries a session-naming
accept hint:

\begin{schema}{MarkerConsistent}
  NotificationState
  \where
  (talkDesc \neq talkDescDefault) \iff \\
  \quad (\dom talkPending \neq \emptyset \lor talkBare \neq \emptyset \lor
  talkQueued > 0 \lor talkConnected = ztrue) \\
  (readDesc \neq readDescDefault) \iff unreadCount > 0
\end{schema}

\begin{schema}{HintNamesSession}
  NotificationState
  \where
  talkDesc \neq talkDescDefault \implies talkBare = \emptyset
\end{schema}

Neither property is yet part of $NotificationState$, so both may fail.
Two model-checking runs at set size~2 exhibit the failures as reachable
states.

\paragraph{The stale marker.}  Setting the ProB goal to
\[
  talkDesc \neq talkDescDefault \land \lnot\, talkActivity
\]
finds a reachable state: an invite arrives ($TalkInviteArrive$ or
$TalkMessageArrive$ raises the marker), then $TalkDrainStale$ empties
every activity source while preserving $talkDesc$.  The marker persists
with nothing behind it.  This is the biff-9la defect.

\paragraph{The bare hint.}  Setting the goal to
\[
  talkDesc \neq talkDescDefault \land talkBare \neq \emptyset
\]
finds a reachable state reached by $TalkInviteArriveBare$: the marker is
present but the pending invite carries no session key, so its accept hint
names no session.

\paragraph{Liveness.}  The safety goals above show a \emph{bad state} is
reachable.  The sharper statement is a liveness one: from some reachable
marker-present state, the marker can never clear.  Expressed as the CTL
formula
\[
  \mathrm{AG}\,(talkDesc \neq talkDescDefault \implies
    \mathrm{EF}\;talkDesc = talkDescDefault),
\]
the model refutes it, because $TalkDrainStale$ is the only way to empty
activity and it does not touch $talkDesc$: once set, the marker is
permanent.  The liveness refutation is robust---it holds even though
other operations remain enabled, so it is not masked the way a plain
deadlock check would be.

%%----------------------------------------------------------------------
\section{The Fix: reconciling transitions}
\label{sec:markerfix}
%%----------------------------------------------------------------------

The repair has two parts.  First, $MarkerConsistent$, $HintNamesSession$
and the read biconditional are promoted into the state schema
(Section~\ref{sec:state}), so they are now invariants that every
operation must respect.  Second, transitions are added that remove
activity \emph{and reconcile the description in the same step}.

Promoting the invariants has an immediate and welcome consequence: the
two defect operations become infeasible.  $TalkDrainStale$ empties
activity while preserving $talkDesc \neq talkDescDefault$; its post-state
now contradicts $MarkerConsistent$, so it is enabled in no reachable
state.  $TalkInviteArriveBare$ sets $talkBare \neq \emptyset$ while
raising the marker; its post-state contradicts $HintNamesSession$, so it
too is dead.  ProB reports both as uncovered operations.  This is the
point of the exercise: once the consistency of marker and activity is an
invariant, the specification \emph{proves} that no drain can strand the
marker and no invite can be recorded without its session key.  The defect
operations are kept in the text only as the negative results they now
are.

A dead operation is not a live transition, however.  Something must
still remove activity, and---because the invariant forbids emptying
activity without reverting the description---any such transition must do
both at once.  Four suffice: $TalkDrain$, $WithdrawArrive$,
$ExpirePendingInvite$ and $TalkAccept$; the proof in
Section~\ref{sec:markerfix2} shows each discharges an obligation the
others do not.  $WithdrawArrive$ and $ExpirePendingInvite$ together
replace the single abstract $DropPendingInvite$ of the earlier drafts:
the two concrete removal paths the implementation must supply---a
withdrawal frame and a time-to-live expiry---have distinct preconditions
and so are kept distinct in the model.

The withdrawal path itself carries a further obligation the earlier drafts
missed, which a second review found.  A withdrawal frame names the invite it
cancels only by the inviter's user identity, but the pending invite is keyed
by session: $talkPending$ records the inviter's $SESSIONKEY$.  If
$WithdrawArrive$ removed on the bare user match, a withdrawal emitted by an
earlier session could---because the talk transport gives no cross-session
ordering---arrive after a fresh invite from a later session of the same user
and silently delete the live invite.  The guard
$talkPending~inviter? = withdrawKey?$ forecloses this: a withdrawal is applied
only when its originating session key matches the stored one.  A frame whose
key does not match is dropped by $WithdrawForeign$, the withdrawal-side
analogue of $talk.tex$'s $DrainForeignAccept$.  The username-only version is
retained as the dead operation $WithdrawStale$, whose reachability under the
missing guard, and unreachability once the guard is present, are the subject
of the reordering proof in Section~\ref{sec:markerfix2}.

\subsection{Draining messages and ending the conversation}

$TalkDrain$ empties the message queue and ends the connection.  If no
pending invite remains behind, it reverts the description; otherwise the
marker stays, now backed only by the surviving invites.

\begin{schema}{TalkDrain}
  \Delta NotificationState
  \where
  pendingAge' = pendingAge \\
  (talkQueued > 0 \lor talkConnected = ztrue) \\
  talkQueued' = 0 \\
  talkConnected' = zfalse \\
  ((\dom talkPending = \emptyset \land talkBare = \emptyset) \implies \\
  \quad (talkDesc' = talkDescDefault \land talkDescChanged' = zfalse)) \\
  ((\dom talkPending \neq \emptyset \lor talkBare \neq \emptyset) \implies \\
  \quad (talkDesc' = talkDesc \land talkDescChanged' = talkDescChanged)) \\
  natsConn' = natsConn \\
  kvProvisioned' = kvProvisioned \\
  talkSubscribed' = talkSubscribed \\
  session' = session \\
  toolContext' = toolContext \\
  wallDesc' = wallDesc \\
  readDesc' = readDesc \\
  wallDescChanged' = wallDescChanged \\
  unreadCount' = unreadCount \\
  talkPending' = talkPending \\
  talkBare' = talkBare \\
  queueItems' = queueItems \\
  queueKinds' = queueKinds \\
  queueTexts' = queueTexts \\
  currentIndex' = currentIndex \\
  fileExists' = fileExists \\
  fileUnreadCount' = fileUnreadCount \\
  fileItems' = fileItems \\
  barLine2' = barLine2 \\
  barStale' = barStale \\
  lastNotifyResult' = lastNotifyResult \\
  pendingNotify' = pendingNotify \\
  notifyLost' = notifyLost \\
  inFlight' = inFlight \\
  beltPnSnapshot' = beltPnSnapshot \\
  capPnSnapshot' = capPnSnapshot \\
  raceLost' = raceLost \\
  staleWithdrawHit' = staleWithdrawHit
\end{schema}

\subsection{Withdrawing a pending invite (ntWithdraw frame)}

The earlier drafts collapsed the two ways a pending invite can be removed
into a single abstract operation.  The two mechanisms have genuinely
different preconditions, however, and the implementation must provide both,
so they are modelled here as distinct operations---the partition of
Section~\ref{sec:markerfix2} then owes a separate test for each.

$WithdrawArrive$ is the first: the recipient receives an \emph{ntWithdraw}
notification from the inviter---the inviter ran \texttt{talk\_end}, hung up,
or otherwise cancelled---and removes that inviter's pending entry.  A
withdrawal frame carries the inviter's originating session key in the field
$withdrawKey?$; this is the \emph{ntWithdraw} analogue of the $nfromKey$ that
$talk.tex$ guards on for accepts.  The removal is honoured only when that key
equals the session key stored against the pending invite---%
$inviter? \in \dom talkPending$ and $talkPending~inviter? = withdrawKey?$.
The guard is not decoration.  The talk transport is core NATS publish and
subscribe; it gives no cross-session ordering, so a withdrawal emitted by an
earlier session of the same user may be delivered \emph{after} a fresh invite
from a later session of that user.  A user invites us from session~$A$,
withdraws, and re-invites from session~$B$; because $talkPending$ is a
function keyed on the user, the re-invite supersedes the first, so
$talkPending$ now maps that user to $B$'s key.  A withdrawal keyed to
session~$A$ that arrives late must not disturb the live session~$B$
invite---and with the key guard it does not, because the stored key (session
$B$'s) differs from the withdrawal's key (session $A$'s), so the equality
$talkPending~inviter? = withdrawKey?$ fails and no removal occurs.  Without
the guard, the withdrawal would delete the current invite on a bare user
match, the marker would clear, and the live invite would be silently lost.  This is
exactly the stale-marker defect class the specification exists to prevent,
reached here by a non-adversarial reordering rather than by an attacker; the
key guard is the withdrawal-side counterpart of the accept-side consent
boundary of $talk.tex$ (DES-043, invariant~11).  A withdrawal whose key does
not match is not applied here; it is handled by $WithdrawForeign$ below.
$WithdrawArrive$'s precondition is otherwise the arrival of the frame on a
live connection; it does \emph{not} depend on how long the invite has been
pending.  If no activity survives, the description reverts and the clock
returns to zero.

\begin{schema}{WithdrawArrive}
  \Delta NotificationState \\
  inviter? : USER \\
  withdrawKey? : SESSIONKEY
  \where
  natsConn = natsConnected \\
  inviter? \in \dom talkPending \\
  talkPending~inviter? = withdrawKey? \\
  talkPending' = \{ inviter? \} \ndres talkPending \\
  talkBare' = talkBare \\
  ((\dom talkPending' = \emptyset \land talkBare' = \emptyset) \implies
  pendingAge' = 0) \\
  ((\dom talkPending' \neq \emptyset \lor talkBare' \neq \emptyset) \implies
  pendingAge' = pendingAge) \\
  ((\dom talkPending' = \emptyset \land talkBare' = \emptyset \land \\
  \quad talkQueued = 0 \land talkConnected = zfalse) \implies \\
  \quad (talkDesc' = talkDescDefault \land talkDescChanged' = zfalse)) \\
  ((\dom talkPending' \neq \emptyset \lor talkBare' \neq \emptyset \lor \\
  \quad talkQueued > 0 \lor talkConnected = ztrue) \implies \\
  \quad (talkDesc' = talkDesc \land talkDescChanged' = talkDescChanged)) \\
  natsConn' = natsConn \\
  kvProvisioned' = kvProvisioned \\
  talkSubscribed' = talkSubscribed \\
  session' = session \\
  toolContext' = toolContext \\
  wallDesc' = wallDesc \\
  readDesc' = readDesc \\
  wallDescChanged' = wallDescChanged \\
  unreadCount' = unreadCount \\
  talkQueued' = talkQueued \\
  talkConnected' = talkConnected \\
  queueItems' = queueItems \\
  queueKinds' = queueKinds \\
  queueTexts' = queueTexts \\
  currentIndex' = currentIndex \\
  fileExists' = fileExists \\
  fileUnreadCount' = fileUnreadCount \\
  fileItems' = fileItems \\
  barLine2' = barLine2 \\
  barStale' = barStale \\
  lastNotifyResult' = lastNotifyResult \\
  pendingNotify' = pendingNotify \\
  notifyLost' = notifyLost \\
  inFlight' = inFlight \\
  beltPnSnapshot' = beltPnSnapshot \\
  capPnSnapshot' = capPnSnapshot \\
  raceLost' = raceLost \\
  staleWithdrawHit' = staleWithdrawHit
\end{schema}

\subsection{A withdrawal that names the wrong session (WithdrawForeign)}

$WithdrawForeign$ is the withdrawal-side counterpart of $talk.tex$'s
$DrainForeignAccept$.  A withdrawal frame arrives whose key does \emph{not}
match the pending invite it purports to cancel---either the user has no
pending invite at all, or the stored invite carries a different session key,
as in the reordering above.  The frame is dropped: the pending set, the
marker, the clock, and every other component are left exactly as they were.
Modelling it as a distinct operation, rather than as a silent no-op, forces
the test-partition of Section~\ref{sec:markerfix2} to owe a case for it, and
records in the specification that the implementation must recognise and
discard such a frame rather than acting on its user field alone.

\begin{schema}{WithdrawForeign}
  \Delta NotificationState \\
  inviter? : USER \\
  withdrawKey? : SESSIONKEY
  \where
  natsConn = natsConnected \\
  \lnot\, (inviter? \in \dom talkPending \land
  talkPending~inviter? = withdrawKey?) \\
  talkPending' = talkPending \\
  talkBare' = talkBare \\
  pendingAge' = pendingAge \\
  talkDesc' = talkDesc \\
  talkDescChanged' = talkDescChanged \\
  natsConn' = natsConn \\
  kvProvisioned' = kvProvisioned \\
  talkSubscribed' = talkSubscribed \\
  session' = session \\
  toolContext' = toolContext \\
  wallDesc' = wallDesc \\
  readDesc' = readDesc \\
  wallDescChanged' = wallDescChanged \\
  unreadCount' = unreadCount \\
  talkQueued' = talkQueued \\
  talkConnected' = talkConnected \\
  queueItems' = queueItems \\
  queueKinds' = queueKinds \\
  queueTexts' = queueTexts \\
  currentIndex' = currentIndex \\
  fileExists' = fileExists \\
  fileUnreadCount' = fileUnreadCount \\
  fileItems' = fileItems \\
  barLine2' = barLine2 \\
  barStale' = barStale \\
  lastNotifyResult' = lastNotifyResult \\
  pendingNotify' = pendingNotify \\
  notifyLost' = notifyLost \\
  inFlight' = inFlight \\
  beltPnSnapshot' = beltPnSnapshot \\
  capPnSnapshot' = capPnSnapshot \\
  raceLost' = raceLost \\
  staleWithdrawHit' = staleWithdrawHit
\end{schema}

\subsection{The withdrawal gap, as a dead operation (WithdrawStale)}

$WithdrawStale$ is the defect this section exists to close, kept in the text
as the negative result it now is.  It models the username-only withdrawal:
an \emph{ntWithdraw} frame arrives whose key does not match the stored invite
($talkPending~inviter? \neq withdrawKey?$), yet it removes the invite anyway,
on the bare user match.  The removed invite is the live one; the marker
clears with nothing behind it and the current invite is silently lost.  The
history component $staleWithdrawHit$ records that this has happened, so the
suppression can be named as a reachable state.

Promoting $staleWithdrawHit = zfalse$ into $NotificationState$
(Section~\ref{sec:state}) makes $WithdrawStale$ infeasible: its post-state
sets $staleWithdrawHit' = ztrue$, which the invariant forbids.  ProB reports
it as an uncovered operation, and the goal $staleWithdrawHit = ztrue$ is
unreachable.  Delete the invariant---as an earlier, username-only draft
implicitly did---and $WithdrawStale$ becomes enabled: with two distinct
session keys it drives the model to $staleWithdrawHit = ztrue$, the reordering
counterexample of Section~\ref{sec:markerfix2}.  This is the same device by
which $TalkDrainStale$ and $TalkInviteArriveBare$ are proved dead against
$MarkerConsistent$ and $HintNamesSession$.

\begin{schema}{WithdrawStale}
  \Delta NotificationState \\
  inviter? : USER \\
  withdrawKey? : SESSIONKEY
  \where
  natsConn = natsConnected \\
  inviter? \in \dom talkPending \\
  talkPending~inviter? \neq withdrawKey? \\
  talkPending' = \{ inviter? \} \ndres talkPending \\
  talkBare' = talkBare \\
  staleWithdrawHit' = ztrue \\
  ((\dom talkPending' = \emptyset \land talkBare' = \emptyset) \implies
  pendingAge' = 0) \\
  ((\dom talkPending' \neq \emptyset \lor talkBare' \neq \emptyset) \implies
  pendingAge' = pendingAge) \\
  ((\dom talkPending' = \emptyset \land talkBare' = \emptyset \land \\
  \quad talkQueued = 0 \land talkConnected = zfalse) \implies \\
  \quad (talkDesc' = talkDescDefault \land talkDescChanged' = zfalse)) \\
  ((\dom talkPending' \neq \emptyset \lor talkBare' \neq \emptyset \lor \\
  \quad talkQueued > 0 \lor talkConnected = ztrue) \implies \\
  \quad (talkDesc' = talkDesc \land talkDescChanged' = talkDescChanged)) \\
  natsConn' = natsConn \\
  kvProvisioned' = kvProvisioned \\
  talkSubscribed' = talkSubscribed \\
  session' = session \\
  toolContext' = toolContext \\
  wallDesc' = wallDesc \\
  readDesc' = readDesc \\
  wallDescChanged' = wallDescChanged \\
  unreadCount' = unreadCount \\
  talkQueued' = talkQueued \\
  talkConnected' = talkConnected \\
  queueItems' = queueItems \\
  queueKinds' = queueKinds \\
  queueTexts' = queueTexts \\
  currentIndex' = currentIndex \\
  fileExists' = fileExists \\
  fileUnreadCount' = fileUnreadCount \\
  fileItems' = fileItems \\
  barLine2' = barLine2 \\
  barStale' = barStale \\
  lastNotifyResult' = lastNotifyResult \\
  pendingNotify' = pendingNotify \\
  notifyLost' = notifyLost
\end{schema}

\subsection{Expiring a pending invite (time-to-live)}

The second mechanism covers the inviter who never sends a withdrawal---a
crash, a dropped connection, an abandoned session.  No frame arrives, so
$WithdrawArrive$ is never enabled for that invite; the only thing that
changes is the clock.  $AgeTick$ models the passage of time on a poller
tick: while any invite is pending and the clock is below its bound, the
clock advances by one.  It touches nothing else---not the pending set, not
the description, not the activity.

\begin{schema}{AgeTick}
  \Delta NotificationState
  \where
  toolContext = zfalse \\
  (\dom talkPending \neq \emptyset \lor talkBare \neq \emptyset) \\
  pendingAge < maxInviteAge \\
  pendingAge' = pendingAge + 1 \\
  natsConn' = natsConn \\
  kvProvisioned' = kvProvisioned \\
  talkSubscribed' = talkSubscribed \\
  session' = session \\
  toolContext' = toolContext \\
  wallDesc' = wallDesc \\
  talkDesc' = talkDesc \\
  readDesc' = readDesc \\
  wallDescChanged' = wallDescChanged \\
  talkDescChanged' = talkDescChanged \\
  unreadCount' = unreadCount \\
  talkPending' = talkPending \\
  talkBare' = talkBare \\
  talkQueued' = talkQueued \\
  talkConnected' = talkConnected \\
  queueItems' = queueItems \\
  queueKinds' = queueKinds \\
  queueTexts' = queueTexts \\
  currentIndex' = currentIndex \\
  fileExists' = fileExists \\
  fileUnreadCount' = fileUnreadCount \\
  fileItems' = fileItems \\
  barLine2' = barLine2 \\
  barStale' = barStale \\
  lastNotifyResult' = lastNotifyResult \\
  pendingNotify' = pendingNotify \\
  notifyLost' = notifyLost \\
  inFlight' = inFlight \\
  beltPnSnapshot' = beltPnSnapshot \\
  capPnSnapshot' = capPnSnapshot \\
  raceLost' = raceLost \\
  staleWithdrawHit' = staleWithdrawHit
\end{schema}

$ExpirePendingInvite$ then reaps an invite once its clock has reached the
bound.  Its precondition---$pendingAge = maxInviteAge$ on a poller
tick---is distinct from $WithdrawArrive$'s, which fires on frame receipt at
any age.  The state effect is the same: the inviter's entry is removed, and
if no activity survives, the description reverts and the clock returns to
zero.

\begin{schema}{ExpirePendingInvite}
  \Delta NotificationState \\
  inviter? : USER
  \where
  toolContext = zfalse \\
  pendingAge = maxInviteAge \\
  (inviter? \in \dom talkPending \lor inviter? \in talkBare) \\
  talkPending' = \{ inviter? \} \ndres talkPending \\
  talkBare' = talkBare \setminus \{ inviter? \} \\
  ((\dom talkPending' = \emptyset \land talkBare' = \emptyset) \implies
  pendingAge' = 0) \\
  ((\dom talkPending' \neq \emptyset \lor talkBare' \neq \emptyset) \implies
  pendingAge' = pendingAge) \\
  ((\dom talkPending' = \emptyset \land talkBare' = \emptyset \land \\
  \quad talkQueued = 0 \land talkConnected = zfalse) \implies \\
  \quad (talkDesc' = talkDescDefault \land talkDescChanged' = zfalse)) \\
  ((\dom talkPending' \neq \emptyset \lor talkBare' \neq \emptyset \lor \\
  \quad talkQueued > 0 \lor talkConnected = ztrue) \implies \\
  \quad (talkDesc' = talkDesc \land talkDescChanged' = talkDescChanged)) \\
  natsConn' = natsConn \\
  kvProvisioned' = kvProvisioned \\
  talkSubscribed' = talkSubscribed \\
  session' = session \\
  toolContext' = toolContext \\
  wallDesc' = wallDesc \\
  readDesc' = readDesc \\
  wallDescChanged' = wallDescChanged \\
  unreadCount' = unreadCount \\
  talkQueued' = talkQueued \\
  talkConnected' = talkConnected \\
  queueItems' = queueItems \\
  queueKinds' = queueKinds \\
  queueTexts' = queueTexts \\
  currentIndex' = currentIndex \\
  fileExists' = fileExists \\
  fileUnreadCount' = fileUnreadCount \\
  fileItems' = fileItems \\
  barLine2' = barLine2 \\
  barStale' = barStale \\
  lastNotifyResult' = lastNotifyResult \\
  pendingNotify' = pendingNotify \\
  notifyLost' = notifyLost \\
  inFlight' = inFlight \\
  beltPnSnapshot' = beltPnSnapshot \\
  capPnSnapshot' = capPnSnapshot \\
  raceLost' = raceLost \\
  staleWithdrawHit' = staleWithdrawHit
\end{schema}

\subsection{Accepting an invite}

$TalkAccept$ is the ordinary path: a pending invite is consumed and the
conversation becomes live.  Activity does not vanish---it moves from
$talkPending$ to $talkConnected$---so the marker is preserved, consistent
with the invariant.

\begin{schema}{TalkAccept}
  \Delta NotificationState \\
  inviter? : USER
  \where
  inviter? \in \dom talkPending \\
  talkPending' = \{ inviter? \} \ndres talkPending \\
  talkConnected' = ztrue \\
  ((\dom talkPending' = \emptyset \land talkBare' = \emptyset) \implies
  pendingAge' = 0) \\
  ((\dom talkPending' \neq \emptyset \lor talkBare' \neq \emptyset) \implies
  pendingAge' = pendingAge) \\
  talkDesc' = talkDesc \\
  talkDescChanged' = talkDescChanged \\
  natsConn' = natsConn \\
  kvProvisioned' = kvProvisioned \\
  talkSubscribed' = talkSubscribed \\
  session' = session \\
  toolContext' = toolContext \\
  wallDesc' = wallDesc \\
  readDesc' = readDesc \\
  wallDescChanged' = wallDescChanged \\
  unreadCount' = unreadCount \\
  talkBare' = talkBare \\
  talkQueued' = talkQueued \\
  queueItems' = queueItems \\
  queueKinds' = queueKinds \\
  queueTexts' = queueTexts \\
  currentIndex' = currentIndex \\
  fileExists' = fileExists \\
  fileUnreadCount' = fileUnreadCount \\
  fileItems' = fileItems \\
  barLine2' = barLine2 \\
  barStale' = barStale \\
  lastNotifyResult' = lastNotifyResult \\
  pendingNotify' = pendingNotify \\
  notifyLost' = notifyLost \\
  inFlight' = inFlight \\
  beltPnSnapshot' = beltPnSnapshot \\
  capPnSnapshot' = capPnSnapshot \\
  raceLost' = raceLost \\
  staleWithdrawHit' = staleWithdrawHit
\end{schema}

%%----------------------------------------------------------------------
\section{Marker Consistency: the proof}
\label{sec:markerfix2}
%%----------------------------------------------------------------------

With the invariants promoted and the reconciling transitions in place, the
checks that discharge the obligations are run at $DEFAULT\_SETSIZE~2$ with
$USER$ shrunk to one and $SESSIONKEY$ kept at two.  Two session keys are
essential here: the reordering defect turns on one user holding invites from
two distinct sessions, and a single session key would collapse the two into
one and hide the fault.  One user suffices, because every marker property is
a statement about the \emph{aggregate} of the activity sources, not about how
many users contribute to it; shrinking $USER$ keeps the state space
exhaustible while losing no distinction the properties depend on.  At this
scope the exploration reports $INITIALISATION$:$8$ against
$SETUP\_CONSTANTS$:$8$---the $MAX\_INITIALISATIONS$ bound never binds---and
terminates with ``all open states visited'': $598{,}824$ states,
$9{,}938{,}416$ transitions, a complete check, not a truncated one.

\begin{verbatim}
probcli docs/notification.tex -model_check -noltl -coverage \
  -goal "staleWithdrawHit = ztrue" \
  -p DEFAULT_SETSIZE 2 \
  -card USER 1 -card SOURCEKEY 1 -card SESSIONKEY 2 \
  -card QSLOT 1 -card USLOT 1 -card ASLOT 1 -card PSLOT 1 \
  -p MAX_INITIALISATIONS 16 -p MAX_OPERATIONS 8000
\end{verbatim}

\paragraph{Safety.}  Ordinary model checking visits every reachable state
and finds no invariant violation and no deadlock.  Because
$MarkerConsistent$ and $HintNamesSession$ are now part of the state
schema, this establishes that no reachable state has a marker without its
activity, and none has a talk marker whose hint names no session.  The
two goal predicates that were reachable in
Section~\ref{sec:markerprops}---$talkDesc \neq talkDescDefault \land
\lnot\, talkActivity$ and $talkDesc \neq talkDescDefault \land talkBare
\neq \emptyset$---are now unreachable, each being exactly the negation of a
promoted invariant.

\paragraph{The reordering defect, reached and then closed.}  The withdrawal
key guard is proved the same way the stale marker was.  Delete the invariant
$staleWithdrawHit = zfalse$ from the state schema, so that the username-only
$WithdrawStale$ becomes feasible, and set the ProB goal to
$staleWithdrawHit = ztrue$.  The check finds a reachable state by an
eleven-step trace whose tail is the reordering exactly:
$TalkInviteArrive(u, k_B)$ records an invite from the user's session~$B$,
then $WithdrawStale(u, k_A)$---a withdrawal carrying the earlier session~$A$'s
key---removes the live invite on the bare user match, clearing the marker with
the current invite silently lost.  Restore the invariant, and the same goal is
searched over the whole reachable space above: it is \emph{not} found, and the
check reports ``all states visited.''  $WithdrawStale$ is then reported among
the uncovered operations, alongside the earlier defect operations
$TalkDrainStale$ and $TalkInviteArriveBare$: the operation that would suppress
a live invite exists in the text but is enabled in no reachable state.  The
key-matched $WithdrawArrive$ and its foreign-drop counterpart $WithdrawForeign$
are both covered.

\paragraph{The pending-set bound.}  This scope shrinks $USER$ to one, which
suffices for every marker property but cannot exercise
$TalkInviteArriveOverflow$: that operation fires only when a \emph{new}
inviter arrives at a full pending set, and with a single user there is no
second inviter to overflow it.  It is therefore reported uncovered here,
alongside the three dead defect operations---but for an entirely different
reason.  It is a live transition awaiting a second user, not a suppressed one.
The exhaustive two-user check of the Global Constants section
($DEFAULT\_SETSIZE~2$ with $PSLOT$ at one, so $maxPending = 1$) covers it:
one user fills the single pending slot, a second arrives and overflows it,
$TalkInviteArriveOverflow$ evicts the incumbent, and the bound
$\# talkPending + \# talkBare \leq maxPending$ holds across all $598{,}825$
reachable states with no invariant violation and no deadlock.  That bound is
the pending-set analogue of the queue cap $maxQueueLen$ of $talk.tex$; it
closes the unbounded-growth denial of service that a wire-controlled inviter
key---each distinct forgery a permanent entry, the time-to-live being eventual
rather than a cap---would otherwise permit.

\paragraph{Liveness.}  The CTL formula
$\mathrm{AG}\,(talkDesc \neq talkDescDefault \implies
\mathrm{EF}\; talkDesc = talkDescDefault)$,
checked at the same scope, is reported \textsc{true}: from every
marker-present state at least one reconciling transition is enabled and
leads, directly or through $TalkAccept$, to a state in which the marker has
cleared.  The key guard does not weaken this---an invite keyed to $k$ can
still be cleared by a matching $WithdrawArrive(u, k)$, and unconditionally by
the time-to-live $ExpirePendingInvite$---so the channel remains
self-clearing.

\paragraph{Necessity of each transition.}  Each reconciling transition
discharges a distinct obligation.  $TalkDrain$ is the only transition
that clears $talkQueued$ or $talkConnected$; without it a marker backed
by a queued message or a live connection can never clear.
$WithdrawArrive$ and $ExpirePendingInvite$ are the two transitions that
remove a pending invite without the recipient's acceptance, and they are
enabled under different conditions: $WithdrawArrive$ on receipt of an
\emph{ntWithdraw} frame \emph{whose key matches the pending invite}, from an
inviter who cancels; $ExpirePendingInvite$ on a poller tick once the clock
has reached $maxInviteAge$.  The second is
what covers exactly the biff-9la scenario---an invite from a peer who
never returns and never withdraws: only the time-to-live path clears it,
so without $ExpirePendingInvite$ (and the $AgeTick$ that drives its
clock) the talk channel would not be self-clearing in the way the read
reference is.  $WithdrawForeign$ removes nothing; it is the transition that
\emph{declines} to act on a mismatched frame, and it is kept as an explicit
operation so that the implementation owes a case for recognising and
discarding such a frame rather than acting on its user field alone.
$TalkAccept$ is the ordinary progress transition; it is not
required for the marker to clear, but without it the only exit from a
pending invite is to remove it, and the handshake modelled in $talk.tex$
could never complete.

%%----------------------------------------------------------------------
\section{System Properties}
%%----------------------------------------------------------------------

\subsection{Notification Delivery Guarantee}

\paragraph{Erratum.}  Earlier drafts stated $NotifyGuarantee$ below as
an unconditional invariant---``no operation produces $notifyDropped$
after a description change.''  $PollTickNotifyFail$ and
$PollTickNotifyFailUnrecorded$ (Section~\ref{sec:sessionloss})
falsify it directly: both are enabled exactly when a description has
changed, and both produce $lastNotifyResult' = notifyDropped$.  The
schema is kept below, unchanged, as a record of what was believed and
is no longer claimed as a property of the whole reachable space.  What
\emph{does} hold in the whole reachable space is the accounting safety
property $notifyLost = zfalse$ (Section~\ref{sec:sessionloss}): a
dropped send is not guaranteed never to happen, but it is guaranteed
never to happen unrecorded.

\begin{schema}{NotifyGuarantee}
  NotificationState
  \where
  (wallDescChanged = ztrue \lor talkDescChanged = ztrue)
  \implies lastNotifyResult = notifySent
\end{schema}

\emph{What remains true, and a modelling seam worth naming}: $WallPost$
and $TalkStart$ still set $lastNotifyResult' = notifySent$
unconditionally.  This is not, on reflection, a claim that the belt
path they use is somehow exempt from $NotifyBeltSendFail$
(Section~\ref{sec:beltpath})---it is a scoping boundary of this
specification.  $WallPost$/$TalkStart$ model the tool handler that
mutates state and notifies \emph{inline}, in one atomic step, as an
abstraction of the whole request; $NotifyBelt$/$NotifyBeltBegin$ model
the underlying \texttt{notify\_tool\_list\_changed()} call in isolation,
as used by $KVWallReceive$ and $NatsTalkCallback$'s downstream poller
tick.  Splitting $WallPost$/$TalkStart$ the same way $NotifyBelt$ was
split is a straightforward, mechanical extension of this section's
technique should a future finding need it; it is out of scope here
because no review has found a defect that requires it.  $KVWallReceive$
and $NatsTalkCallback$ still defer the actual send to the poller, which
is exactly the design that makes $PollTick$'s split
(Section~\ref{sec:sessionloss}) the right place to model the
suspenders-side failure.

\subsection{Session Invariant---retracted}

\paragraph{Erratum.}  Earlier drafts asserted $SessionInvariant$ below
as a state invariant and used it to prove $SuspendersAvailable$.
Section~\ref{sec:state}'s erratum retracts it: $PollTickNotifyFail$
sets $session' = sessNone$ while $natsConn' = natsConnected$, so the
schema below is false of the corrected model and is not part of it.
The property it was protecting---$EnsureConnected$ cannot run before a
first capture---is now $EnsureConnected$'s own explicit precondition
(Section~\ref{sec:ensureconnected}).

\begin{schema}{SessionInvariant}
  NotificationState
  \where
  natsConn = natsConnected \implies session = sessCaptured
\end{schema}

\subsection{Suspenders Always Available---retracted}

\paragraph{Erratum.}  For the same reason, $SuspendersAvailable$ below
no longer holds unconditionally: after $PollTickNotifyFail$,
$natsConn = natsConnected \land toolContext = zfalse$ while
$session = sessNone$, so the suspenders path is exactly \emph{not}
available until $CaptureSession$, $RecaptureSession$,
$CaptureSessionFlushOk$, or a successful $NotifyBelt$ call restores it.
This is the substance of root cause~\#1: the property below was the
promise the background poller made to itself, and it is the promise
that broke.

\begin{schema}{SuspendersAvailable}
  NotificationState
  \where
  natsConn = natsConnected \\
  toolContext = zfalse \\
  session = sessCaptured
\end{schema}

What holds instead: whenever $SuspendersAvailable$ fails with a changed
flag set, $notifyLost = zfalse$ (the corrected model's safety property)
guarantees $pendingNotify = ztrue$, so the debt is visible to whichever
of $CaptureSession$, $RecaptureSession$, or $CaptureSessionFlushOk$
arrives next---and, once Section~\ref{sec:race}'s mutex is in place,
$raceLost = zfalse$ guarantees that recovery is not itself clobbered by
a second site's concurrent, stale write.

\subsection{Wall and Talk Symmetry}

Wall posting uses the belt path (tool context) via $WallPost$, which
sets $lastNotifyResult' = notifySent$ unconditionally within this
specification's chosen scope (the modelling seam noted above the
$NotifyGuarantee$ erratum).  Wall reception from another user uses
$KVWallReceive$, and talk reception uses $NatsTalkCallback$; both defer
the actual send to the poller rather than producing $notifySent$
themselves, and that deferred send can now fail
(Section~\ref{sec:sessionloss}).  Section~\ref{sec:sessionloss}'s
erratum to $SuspendersAvailable$ applies to both deferred paths: the
poller's next tick attempts the send, and $notifyLost = zfalse$ together
with $raceLost = zfalse$, not delivery symmetry, is what guarantees the
attempt is never silently and unrecoverably lost.

\paragraph{Erratum (biff-9la).}  The symmetry stated above concerns only
\emph{notification delivery}---that a \texttt{tools/list\_changed} is
sent.  It does not extend to the \emph{life-cycle of the tool
description}, and the earlier drafts of this specification wrongly implied
that it did by modelling $talkDesc$ as an opaque value with no invariant.
Wall and talk are not symmetric here.  The wall and read markers are
driven by sources that decrease---an unread count that falls to zero, a
display item that is advanced away---so they clear themselves.  The talk
marker is driven by $talkPending$, $talkQueued$ and $talkConnected$, and
nothing in the original model reverted $talkDesc$ when those emptied:
Sections~\ref{sec:markerprops} reproduces the resulting stale marker as a
reachable state.  The corrected treatment requires the reconciling
transitions of Section~\ref{sec:markerfix}; only then does the talk
description track its activity as the read reference does.  The
consistency invariant $MarkerConsistent$, not delivery symmetry, is the
property that matters.

\paragraph{Erratum ($talk.tex$ invariant~4).}  Invariant~4 of $talk.tex$
(``Pending invites \ldots remain in the set until the next \texttt{talk}
command'') describes exactly the grow-only behaviour that strands the
marker.  If pending invites are never removed except by a subsequent
\texttt{talk} command, then a marker driven by $talkPending$ can never
clear on its own.  Under the invariant $MarkerConsistent$ that reading is
superseded: a pending invite must also be removable by $WithdrawArrive$---an
\emph{ntWithdraw} frame from the inviter \emph{whose session key matches the
stored invite}---or by $ExpirePendingInvite$---a time-to-live expiry aged out
on a poller tick---so that the talk channel is self-clearing.  The key
qualification on $WithdrawArrive$ is not incidental: a withdrawal matched on
the inviter's user identity alone would, under the transport's lack of
cross-session ordering, let a stale withdrawal from one session delete a live
invite from another, reintroducing the stranded-marker defect by a
non-adversarial path.  $WithdrawArrive$ therefore guards on
$talkPending~inviter? = withdrawKey?$, and a mismatched frame is dropped by
$WithdrawForeign$; this is the withdrawal-side mirror of the accept-side
consent guard $(head~queue).nfromKey = partnerKey$ that $talk.tex$'s
$DrainAcceptWhileInviting$ already imposes (DES-043).  The consent-boundary
guarantees of $talk.tex$ (invariant~11) are unaffected: withdrawal and expiry
only \emph{remove} pending invites, never complete a handshake, so they cannot
forge consent.

$WallNotifySuccess$ and $TalkNotifySuccess$ describe the state
immediately after a successful send---$PollTickNotifyOk$, or $NotifyBelt$
before it clears the changed flag.  Neither is a state invariant: a
changed flag can also coexist with $lastNotifyResult = notifyDropped$,
which is exactly $PollTickNotifyFail$'s and
$PollTickNotifyFailUnrecorded$'s post-state
(Section~\ref{sec:sessionloss}).

\begin{schema}{WallNotifySuccess}
  NotificationState
  \where
  wallDescChanged = ztrue \\
  lastNotifyResult = notifySent
\end{schema}

\begin{schema}{TalkNotifySuccess}
  NotificationState
  \where
  talkDescChanged = ztrue \\
  lastNotifyResult = notifySent
\end{schema}

\subsection{Startup Sequence}

The correct startup sequence is:
\begin{enumerate}
  \item $Init$ --- all state zeroed, NATS disconnected, session uncaptured.
  \item $CaptureSession$ --- session captured from server lifespan context.
  \item $EnsureConnected$ --- NATS connected ($EnsureConnected$'s own
    precondition, Section~\ref{sec:ensureconnected}, satisfied).
  \item Poller starts, subscriptions begin --- suspenders path available,
    until and unless a send failure strands it
    (Section~\ref{sec:sessionloss}).
\end{enumerate}

$EnsureConnected$ before any $CaptureSession$ is rejected by its own
guard $session \neq sessNone \lor natsConn = natsConnected$
(Section~\ref{sec:ensureconnected}): the first disjunct fails because
no capture has happened, and the second fails because nothing has
connected yet.  A \emph{mid-session} loss of $session$, by contrast,
does not block a repeat $EnsureConnected$---the second disjunct now
holds, correctly, since the connection itself was never the problem.

\end{document}
